% SPDX-FileCopyrightText: 2026 Will Cook
% SPDX-License-Identifier: CC-BY-4.0
%
% Standalone Erdős Problem Note for #257.  Context is deliberately repeated
% from the gateway so that this file remains independently readable.
%
% Round-5 editorial ordering: theorem, displacement mechanism, sufficient
% extensions and their mathematical boundaries, then the actual target.
\documentclass[11pt]{article}
% SPDX-FileCopyrightText: 2026 Will Cook
% SPDX-License-Identifier: CC-BY-4.0
%
% Shared preamble for the Erdős Problem Notes: the short, standalone,
% problem-owned manuscripts covering the expansion library `ErdosProblems`.
%
% One note owns one Erdős problem.  Each note repeats whatever context its own
% reader needs, so that a reader who arrives at a single problem never has to
% assemble it from the gateway paper.  Deliberate repetition across notes is the
% design, not an oversight.
%
% Source links resolve against one pinned formal-source commit, declared once
% below.  `scripts/check_problem_note_sources.py` verifies that every authored
% (file, line, declaration) triple names that exact declaration in the pinned
% snapshot, so a link cannot silently rot as later work moves lines.

\usepackage{paper-house-style}
\usepackage{array}
\usepackage{booktabs}
\usepackage{enumitem}
\setlist{topsep=0.35em,itemsep=0.2em}
\usepackage[colorlinks=true,linkcolor=linkinternal,urlcolor=linkexternal,
  citecolor=linkinternal]{hyperref}
\hypersetup{bookmarksnumbered=true,bookmarksopen=true,bookmarksopenlevel=1,
  pdfauthor={Will Cook},pdflang={en-GB}}

% ---------------------------------------------------------------------------
% Pinned formal source.  \commit names the checkpoint every link resolves
% against; the checker reads that snapshot rather than the working tree, so the
% printed coordinates stay correct however later waves move declarations.
% ---------------------------------------------------------------------------
\newcommand{\commit}{e5fd26a6d1b1a346430205eab5385b4c9565d004}
\newcommand{\commitshort}{e5fd26a6d1b1}
\newcommand{\repobase}{https://github.com/wcook04/plectis-lean-erdos249-257/blob/\commit}
\newcommand{\sourceurl}{https://github.com/wcook04/plectis-lean-erdos249-257/tree/\commit}
\newcommand{\PX}{ErdosProblems}

% Declaration names stay inside link targets.  The printed label is either a
% spaced, lower-case reading of the declaration name (\lref) or a short authored
% phrase (\lword); neither prints a module path or a raw Lean identifier.
\ExplSyntaxOn
\cs_new_protected:Npn \noteleanlabel #1
  {
    \tl_set:Nx \l_tmpa_tl { \tl_to_str:n {#1} }
    \regex_replace_all:nnN { _+ } { \c{space} } \l_tmpa_tl
    \regex_replace_all:nnN
      { ([a-z0-9])([A-Z]) }
      { \1\c{space}\2 }
      \l_tmpa_tl
    \tl_set:Nx \l_tmpa_tl { \tl_lower_case:n { \tl_use:N \l_tmpa_tl } }
    \tl_use:N \l_tmpa_tl
  }
\ExplSyntaxOff
\newcommand{\leanlabel}[1]{\noteleanlabel{#1}}
\newcommand{\lref}[3]{\href{\repobase/\PX/#1\#L#2}{\leanlabel{#3}}}
\newcommand{\lword}[4]{\href{\repobase/\PX/#1\#L#2}{#4}}
\newcommand{\lloc}[2]{\href{\repobase/\PX/#1\#L#2}{Lean source}}

% Links into a sibling library, named repository-relative.  The expansion
% library is implicit in \lref/\lword, because that is where problem-owned
% work lands.  Several problems have their headline theorem in the older
% reviewed #249/#257 corpus, and a note that could not name that library
% would have to paraphrase a checked statement instead of linking it.
\newcommand{\mref}[3]{\href{\repobase/#1\#L#2}{\leanlabel{#3}}}
\newcommand{\mword}[4]{\href{\repobase/#1\#L#2}{#4}}
\newcommand{\mloc}[2]{\href{\repobase/#1\#L#2}{Lean source}}
\emergencystretch=2em

\newtheorem{theorem}{Theorem}[section]
\newtheorem{proposition}[theorem]{Proposition}
\newtheorem{lemma}[theorem]{Lemma}
\newtheorem{corollary}[theorem]{Corollary}
\theoremstyle{definition}
\newtheorem{definition}[theorem]{Definition}
\newtheorem{example}[theorem]{Example}
\newtheorem{problem}[theorem]{Problem}
\theoremstyle{remark}
\newtheorem*{remark}{Remark}

\newcommand{\N}{\mathbb{N}}
\newcommand{\Npos}{\mathbb{N}_{>0}}
\newcommand{\Z}{\mathbb{Z}}
\newcommand{\Q}{\mathbb{Q}}
\newcommand{\R}{\mathbb{R}}
\newcommand{\lcm}{\operatorname{lcm}}
\newcommand{\ph}{\varphi}

% The series line printed under every note's title block.
\newcommand{\seriesline}{Erd\H{o}s Problem Notes}

% Production provenance.  The wording is fixed by the corpus provenance
% contract and reproduced verbatim in every manuscript in this repository, so
% the papers cannot drift into different accounts of how they were made.
\newcommand{\noteprovenance}{\provenancenote{\emph{Authorship and AI use.} The
prose, formal proofs, and software in this version were drafted and revised by
large-language-model agents under Will Cook's direction.  Cook set the
objectives and acceptance criteria, reviewed and authorised the manuscript, and
is responsible for its contents.  The agents are production tools, not authors.
See \emph{Statements and declarations}.}}

% The status paragraph every note carries, in its own words, before any result.
% Kernel checking establishes the exact Lean proposition; it does not establish
% that the proposition is the interesting one, that it is new, or that the
% Erdős problem is closed.
\newcommand{\statusboundary}{%
  \emph{Status.}  The problem treated here is open, and this note does not
  close it.  Every statement below marked as checked is a proposition that the
  pinned Lean kernel accepts from the sources this note links to, with no
  \texttt{sorry}, no added axiom, and no unchecked evaluation.  That is a claim
  about the formal statement, not about its mathematical interest, its
  novelty, or the original problem.  The unresolved obligations are named
  exactly, in their own section, and none of the finite computations,
  reductions, or no-go results here removes one of them.}

\renewcommand{\commit}{99f4bf47422abbd8757cbb22b50ba079d764d3a7}
\renewcommand{\commitshort}{99f4bf47422a}
\hypersetup{
  pdftitle={Weighted Support Criteria for Reciprocal Mersenne Subseries (Erdős Problem 257)},
  pdfsubject={Erdős Problem Notes: Problem 257},
  pdfkeywords={irrationality, Mersenne numbers, Lambert series, divisor incidence, scaled tails, achievement sets, squarefree, Lean 4},
}

\newcommand{\Ach}{\mathcal A}

\runninghead{Erd\H{o}s \#257: weighted support criteria}
\title{Weighted Support Criteria for Reciprocal Mersenne Subseries}
\subtitle{The series $\sum_{a\in A}(b^a-1)^{-1}$ and Erd\H{o}s Problem~\#257}
\author{Will Cook}
\date{18 September 2026}

\begin{document}
\maketitle
\noteprovenance
\begin{abstract}
For a finite nonempty set $P$ of primes, let $h(a)$ be the $P$-part of $a$.
We prove that $\sum_{a\in A}h(a)/[a(2^{h(a)}-1)]<\infty$ makes
$\sum_{a\in B}(b^a-1)^{-1}$ irrational for every integer $b\ge2$ and every
infinite $B\subseteq A$.  The condition permits a divergent reciprocal
sum, but excludes full support and all odd exponents.  The proof averages
positive displacements over multiples of a finite modulus; a second,
dyadic average controls incomplete periods.  We also prove the
reciprocal-summable criterion stated by Erd\H{o}s and a common-average
extension using positive divisor covers.  Supplementary results concern
finite denominators and conditional rational-membership tests, not a
resolution of the universal problem.

\end{abstract}

\section{Introduction and main results}\label{sec:problem}
For $A\subseteq\Npos$ and $b>1$, put
$X_A(b)=\sum_{a\in A}(b^a-1)^{-1}$.  This converges because
$(b^a-1)^{-1}\le b^{1-a}/(b-1)$.  For a finite set $P$ of primes,
the \emph{$P$-part} of $a$ is $h(a)=\prod_{p\in P}p^{v_p(a)}$, where
$v_p(a)$ is the exponent of $p$ in $a$.  For example, if $P=\{2\}$
and $a=2^km$ with $m$ odd, then $h(a)=2^k$.

\begin{theorem}[a weighted condition on the support]
\label{res:weighted-support}
Let $b\ge2$ be an integer, let $A\subseteq\Npos$ be infinite, and let $P$
be a finite nonempty set of primes.  Set
$h(a)=\prod_{p\in P}p^{v_p(a)}$.  If
\begin{equation}\label{eq:weighted-fixed-base}
 W_{b,P}(A):=\sum_{a\in A}
 \frac{h(a)}{a(b^{h(a)}-1)}<\infty,
\end{equation}
then $X_A(b)$ is irrational.  In particular, the base-two condition
\begin{equation*}\label{eq:weighted-return}
 \sum_{a\in A}\frac{h(a)}{a(2^{h(a)}-1)}<\infty
 \tag{W}
\end{equation*}
implies that $X_A(b)$ is irrational for every integer $b\ge2$.  Both
conclusions are hereditary under passage to infinite subsets.
\end{theorem}

For $P=\{2\}$, the summand in~\eqref{eq:weighted-fixed-base} is
$1/[m(b^{2^k}-1)]$ when $a=2^km$ and $m$ is odd.  Large powers of $2$
can therefore compensate for large reciprocal mass among the odd
factors.  The example following the proof in
Section~\ref{sec:eight-return-extensions} makes this precise.
The reciprocal mass at each fixed $P$-part must be finite, and its
weighted sum over the parts must converge; the first requirement alone
is insufficient.  Section~1.2 of the companion, \emph{Reciprocal Mersenne
Subseries}, gives the decomposition and a separating example.

For every finite $P$, full support and all odd exponents fail the
condition: the integers coprime to $2\prod_{p\in P}p$ have $h(a)=1$
and divergent reciprocal sum, by inclusion--exclusion.  The full set
of primes fails too, since the primes outside $P$ contribute
$1/[a(b-1)]$.  These are limitations of this criterion, not assertions
of rationality.  By contrast, $h/(2^h-1)\le1$ shows that it includes
every reciprocal-summable support.

\begin{theorem}[reciprocal-summable supports]\label{res:reciprocal-support}
Let $A\subseteq\Npos$ be infinite.  If
\[
  \sum_{a\in A}\frac1a<\infty,
\]
then \(X_A(b)\) is irrational for every integer \(b\ge2\).
\end{theorem}

Erd\H{o}s stated Theorem~\ref{res:reciprocal-support}, including the
all-base conclusion, after proving its pairwise-coprime case
\cite[p.~222]{erdos1968}.  We give a complete averaging proof of that stated extension.  Erd\H{o}s also discussed weakening reciprocal
summability~\cite[pp.~222, 226]{erdos1968}; no identification of the
weighted condition with his suggested conditions is asserted.

The proof uses a positive displacement that rationality would separate
from zero.  Erd\H{o}s likewise points to fractional parts of powers times
the series value~\cite[p.~226]{erdos1968}.
For an integer $b\ge2$, infinite $A$ and $N>0$, division of $N$ by
each $a\in A$ gives
\[
 0<\Delta_{b,A}(N)
 :=\sum_{a\in A}\frac{b^{N\bmod a}-1}{b^a-1}
 =(b^N-1)X_A(b)-J_{b,A}(N),\qquad J_{b,A}(N)\in\Z.
 \label{eq:intro-displacement}\tag{D}
\]
The integer term is explicitly
\[
 J_{b,A}(N)=\sum_{\substack{a\in A\\a\le N}}
             \sum_{j=1}^{\lfloor N/a\rfloor}b^{N-ja}.
\]
The identity holds for every real $b>1$; integrality uses the integer-base
hypothesis.  Every summand in the displacement
is nonnegative, and an exponent $a>N$ gives a positive summand; such an
exponent exists because $A$ is infinite. Rationality $X_A(b)=p/q$
would force every positive displacement to be at least $1/q$.
The proof makes these positive displacements arbitrarily small by
averaging along multiples of a growing divisibility modulus.

For an application of Theorem~\ref{res:reciprocal-support}, call an integer
\emph{powerful} when every prime dividing it occurs to exponent at least two.  Every such integer can be written
as $u^2v^3$, so
\[
 \sum_{\substack{a\ge1\\a\text{ powerful}}}\frac1a
 \le\zeta(2)\zeta(3)<\infty.
\]
Thus arbitrary infinite thinnings of the powerful integers are included.
For full perfect-power supports, compare Duverney--Tachiya
\cite[Corollary~1.2, p.~4]{duverneytachiya} and the earlier
linear-independence results of Luca--Tachiya
\cite{lucatachiya2014independence}.

\begin{problem}[Erd\H{o}s \#257]\label{res:problem}
Is $X_A(2)$ irrational for every infinite $A\subseteq\Npos$?
\end{problem}
Sections~\ref{sec:reciprocal-support}--\ref{sec:eight-return-extensions}
prove the support criteria: first the simpler reciprocal-summable
argument, then the weighted theorem and its common-average extension.
The remaining sections give supplementary finite-denominator and
rational-membership results.  None is needed for the weighted proof.

The coefficient sequence throughout is the divisor transform of the
indicator function of the support:
\begin{equation}
 c_A(n)=\#\{a\in A:a\mid n\},\qquad
 X_A(b)=\sum_{n\ge1}\frac{c_A(n)}{b^n}.
 \label{eq:incidence}
\end{equation}
Tonelli's theorem applied to
$(b^a-1)^{-1}=\sum_{j\ge1}b^{-aj}$ proves the identity.  Moreover
$0\le c_A(n)\le\tau(n)$.  The selector $1_A$ is Boolean; its divisor transform $c_A$ generally is not.

\section{Reciprocal-summable supports at every integer base}
\label{sec:reciprocal-support}

A shift divisible by each exponent in a finite part of $A$ makes those
terms of the displacement vanish.  We then average over such shifts.
Reciprocal summability controls the terms outside that finite part without
requiring a common period for the whole support.

Fix the integer base $b\ge2$ and put
\[
 w_{b,d}(N)=\frac{b^{N\bmod d}}{b^d-1},\qquad
 T_N^{(b)}=\sum_{d\in A}w_{b,d}(N).
\]
Then $T_N^{(b)}-T_0^{(b)}=\Delta_{b,A}(N)$.
For a fixed positive integer $Q$, the orbit of $Q$ modulo $d$ consists of
$d/g$ residues, where $g=(Q,d)$.  Summing the geometric progression gives
\begin{equation}
 M_{Q,b}(d):=\lim_{X\to\infty}\frac1X\sum_{m=1}^Xw_{b,d}(Qm)
 =\frac{g}{d(b^g-1)}\le\frac1d.
 \label{eq:gcd-orbit-mean}
\end{equation}
For example, at $b=2$, $Q=4$ and $d=6$, the orbit is $4,2,0$ modulo $6$.
Its mean atom is $(16+4+1)/(3\cdot63)=1/9$, whereas its zero-shift atom is
$1/63$. If $6\mid Q$, the orbit instead stays at zero and the mean is
exactly $1/63$. Making an exponent divide $Q$ therefore changes its orbit mean.  The
infinite tail still requires the bound that follows.
The passage from individual atoms to their infinite sum needs a uniform
bound.  Counting positive multiples of $d$ gives
\begin{align*}
 \frac1X\sum_{m=1}^Xw_{b,d}(Qm)
 &=\sum_{r\ge1}b^{-r}
       \frac{\#\{1\le m\le X:d\mid Qm+r\}}X\\
 &\le\sum_{r\ge1}b^{-r}\frac{Q+r}{d}
 =\frac{Q/(b-1)+b/(b-1)^2}{d}\le\frac{Q+2}{d}.
\end{align*}
Indeed, the counted multiples are distinct and at most $QX+r$.
For this fixed $Q$, the majorant $(Q+2)/d$ is summable over $A$.
Dominated convergence therefore takes the observation-length limit inside
the sum over $d$, giving the Ces\`aro mean
$\sum_{d\in A}M_{Q,b}(d)$.  This majorant is not uniform in growing $Q$.

Next let $Q_t=\operatorname{lcm}(1,\ldots,t)$.
For each fixed $d$, eventually $d\mid Q_t$ and
$M_{Q_t,b}(d)=(b^d-1)^{-1}$.  A second dominated-convergence argument,
using $M_{Q_t,b}(d)\le1/d$, yields
\begin{equation}
 \sum_{d\in A}M_{Q_t,b}(d)\longrightarrow X_A(b).
 \label{eq:lcm-prefix-orbit-limit}
\end{equation}
Thus the limiting averages of the nonnegative displacements
$\Delta_{b,A}(Q_tm)$ tend to zero.  Choose $t$, then a sufficiently long
finite average, and finally a term no larger than that average.
This gives arbitrarily small positive displacements, contradicting the
rational lattice in~\eqref{eq:intro-displacement}.  This proves
Theorem~\ref{res:reciprocal-support} directly at every integer base.

The order of limits is essential: the observation length tends to infinity
with the modulus fixed, and only then does the modulus increase.
Reciprocal summability controls both interchanges.  The exact all-base result,
with hypotheses $b\ge2$, infinitude of $A$, and summability of $a\mapsto
\mathbf1_A(a)/a$, is kernel-checked as
\href{https://github.com/wcook04/plectis-erdos/blob/065e0952d7a3a0c9f87275323d50aaedb78f481f/lean/Erdos249257/AllBaseReciprocalSupportIrrationality.lean#L395}{\texttt{irrational\_erdosSupportSeries\_of\_summable\_reciprocal}}.

\section{Extensions beyond reciprocal summability}
\label{sec:eight-return-extensions}

The weighted condition can hold when $\sum_{a\in A}1/a$ diverges, so the
previous proof's summable majorant is no longer available.  We instead
bound finite averages for each exponent and then average over dyadic
observation lengths.  This second average controls the incomplete periods.
Estimate (S) will also allow the divisor-cover argument to use these same
observation lengths.

\subsection{Proof of the weighted criterion}
For a finite nonempty prime set $P$, put
$h(a)=\prod_{p\in P}p^{v_p(a)}$.

We prove the fixed-base assertion of Theorem~\ref{res:weighted-support}.
The base-two condition implies every integer-base instance because
$b^{h(a)}-1\ge2^{h(a)}-1$.


\begin{proof}
The rational-lattice obstruction is~\eqref{eq:intro-displacement}: if
$X_A(b)=p/q$, every positive $\Delta_{b,A}(m)$ is at least $1/q$.
For $d_a(m)=(b^{m\bmod a}-1)/(b^a-1)$ and $g=(Q,a)$, one complete orbit of
$Qm\bmod a$ gives, for $Q,a,T\ge1$,
\begin{equation}\label{eq:weighted-finite-orbit}
 \frac1T\sum_{t=1}^T d_a(tQ)
 \le \frac{g}{a(b^g-1)}+\frac1{T(b^g-1)}.
\end{equation}
Indeed, the orbit has length $a/g$ and
$\sum_{t=1}^{a/g}b^{tQ\bmod a}/(b^a-1)=1/(b^g-1)$; an incomplete orbit costs
at most one more complete orbit.  Also, for $Y=QT$,
\begin{equation}\label{eq:weighted-outer-short}
 \frac1T\sum_{t=1}^T\sum_{\substack{a\in A\\a>Y}}d_a(tQ)
 \le\frac4T,
\end{equation}
because $d_a(tQ)\le2\,2^{tQ-a}$ when $a>QT$ and
$\sum_{t=1}^T2^{tQ-QT}\le2$.

The incomplete-orbit errors are controlled by a second finite average.  For integers $Q,M\ge1$ and $\alpha_a\ge0$ with $\sum_a\alpha_a/a<\infty$,
\begin{equation}\label{eq:weighted-dyadic-short}
 \sum_{j=M}^{2M-1}\frac1{2^j}\sum_{a\le Q2^j}\alpha_a
 \le2Q\sum_{a\ge1}\frac{\alpha_a}{a};
\end{equation}
for each fixed $a$, the admissible geometric tail is at most $2Q/a$.

Fix $\varepsilon>0$.  Choose finite nonempty $F\subseteq A$ so that the
\eqref{eq:weighted-fixed-base} mass outside $F$ is below $\varepsilon$, and
let $L$ be any fixed positive common multiple of $F$.  With $p_*=\max P$, $r=|P|$, and large
$H\ge2p_*$, set
\[
 Q=L\prod_{p\in P}p^{\lfloor\log_pH\rfloor},\qquad
 G=\left\lfloor\frac H{p_*}\right\rfloor.
\]
For $a\in F$ the displacement term vanishes, since $a\mid Q$.
For $a\notin F$ with $h(a)\le H$, we have $h(a)\mid Q$ and hence
$(Q,a)\ge h(a)$.  The complete-period term
in~\eqref{eq:weighted-finite-orbit} is therefore bounded by
$h(a)/[a(b^{h(a)}-1)]$, using the monotonicity of $n/(b^n-1)$.
Their sum is less than $\varepsilon$; their incomplete-period errors will
be handled by the dyadic average.

For $h(a)>H$, we instead have $(Q,a)\ge G$.  Sum the complete-period
bounds over $a\le QT$, using $\sum_{a\le QT}1/a\le1+\log(QT)$.
There are at most $QT$ incomplete-period terms, each at most
$1/[T(b^G-1)]$.  The total contribution from these exponents is thus at most
\[
 \frac{G(1+\log(QT))+Q}{b^G-1}.
\]
For the gcd claim, either every $P$-prime-power component of $h(a)$ is at
most $H$, in which case $h(a)\mid Q$, or some $p^{v_p(a)}>H$ contributes
$p^{\lfloor\log_pH\rfloor}>H/p\ge H/p_*\ge G$ to the gcd.
The second averaging length must make both $Q/M$ and $GM/b^G$
small.  Since $Q$ grows only polynomially in $H$ and $G$ grows linearly,
$M=\lfloor b^{G/2}\rfloor$ lies between these two scales.
Combining the preceding estimates with~\eqref{eq:weighted-outer-short},
averaging over $T=2^j$ for $M\le j<2M$, and
using~\eqref{eq:weighted-dyadic-short} with
$\alpha_a={\bf1}_A(a)/(b^{h(a)}-1)$ yields
\begin{equation}\label{eq:weighted-main-bound}
 \frac1M\sum_{j=M}^{2M-1}\frac1{2^j}
 \sum_{t=1}^{2^j}\Delta_{b,A}(tQ)
 \le \varepsilon+\frac{2QW_{b,P}(A)}M
 +\frac{G(1+\log Q+2M\log2)+Q}{b^G-1}+4\,2^{-M}.
\end{equation}
The order of choices is important.  First fix $\varepsilon$ and choose
$F$ and its common multiple $L$; only then let $H\to\infty$.  With $L$ fixed,
$Q\le LH^r$, $G=H/p_*+O(1)$ and $M\asymp b^{G/2}$, so every term after
$\varepsilon$ tends to zero.  This produces arbitrarily small displacements, but gives no specified
decay rate in the index $N=Qt$.  The left side is a finite average of positive
displacements, so one is below $2\varepsilon$ for large $H$.
Since $\varepsilon$ is arbitrary, this contradicts the lower bound $1/q$
under rationality.  Finally
$W_{b,P}(A)\le W_{2,P}(A)$ for $b\ge2$, and weighted mass decreases on taking
subsets.
\end{proof}

The exact interface of Theorem~\ref{res:weighted-support}, including its
fixed-base conclusion and hereditary all-base clause, is kernel-checked as
\href{https://github.com/wcook04/plectis-erdos/blob/065e0952d7a3a0c9f87275323d50aaedb78f481f/lean/ErdosProblems/Erdos257/PaperCompleteR8/WeightedReturn.lean#L99}{\texttt{divisibilityWeightedClaim}}.

\subsection{An example beyond reciprocal summability}
The hypothesis is genuinely weaker than reciprocal summability.  Let
\[
 A_\star=\{2^k m:k\ge1,\ m\text{ odd},\ m\le2^{2^k}\}.
\]
Write $H_N=\sum_{m=1}^N1/m$ and $N_k=2^{2^k}$.  The reciprocal mass
in the $k$th layer is
\[
 \rho_k=2^{-k}\sum_{\substack{m\le N_k\\m\text{ odd}}}\frac1m
 =2^{-k}\bigl(H_{N_k}-\tfrac12H_{N_k/2}\bigr)
 \longrightarrow\frac{\log2}{2},
\]
so $\sum_{a\in A_\star}1/a$ diverges.  For $P=\{2\}$ the weighted mass of
that layer is
\[
 \frac{2^k}{2^{2^k}-1}\rho_k.
\]
Since $H_N\le1+\log N$, one has $\rho_k\le2^{-k}+\log2$.
The weighted terms are therefore at most
$(1+2^k\log2)2^{1-2^k}$, a summable sequence.
Thus every infinite subset of $A_\star$ has irrational $X_A(b)$ at
every integer base, although $\sum_{a\in A_\star}1/a$ diverges.

Both proofs select a term no larger than a finite average.  For a
related selection step on arithmetic progressions, see
Duverney--Tachiya~\cite[Section~2, (2.3)--(2.9)]{duverneytachiya}.
The sparse-coefficient criteria of Kaneko--Suzuki--Tachiya
\cite[Theorems~1 and~3]{kanekosuzukitachiya} do not apply directly to
$c_A$: for nonempty $A$, it is positive on every multiple of $\min A$.
The companion, Section~1.2, gives the counting argument and distinguishes
their remote-tail average from the displacement used here.

\subsection{A common finite average for the extensions}\label{sec:common-kernel}
To combine the weighted criterion with positive divisor majorants, both
arguments must use the same indices.  The following estimate supplies
that common average; Theorem~\ref{res:mixed-supports} gives the combination.

For $1<B\le2$, positive integers $L,d,M$, and an integer $R\ge0$, put
\[
 w_{B,d}(n)=\frac{B^{n\bmod d}}{B^d-1},\qquad
 \mathscr D_{L;R,M}F
 =\frac1M\sum_{j=R}^{R+M-1}\frac1{2^j}
      \sum_{m=1}^{2^j}F(Lm).
\]
The finite estimate
\begin{equation*}
 \mathscr D_{L;R,M}w_{B,d}
 \le\frac{1+4L/M}{d(B-1)}
 \tag{S}\label{eq:mixed-finite-kernel}
\end{equation*}
equivalently bounds $(B-1)\mathscr D_{L;R,M}w_{B,d}$ by
$(1+4L/M)/d$.  This normalized bound is uniform as $B\downarrow1$;
the unnormalized right side grows like $(B-1)^{-1}$.
The ratio $L/M$ measures the cost of incomplete modular periods.

To prove it, fix $T=2^j$ and average over $1\le m\le T$.
When $d\le LT$, the complete orbit has length $d/g$, with $g=(L,d)$,
and total weight $1/(B^g-1)$.  Complete cycles and one remaining piece give
\[
 \frac1T\sum_{m=1}^T w_{B,d}(Lm)
 \le\frac{g}{d(B^g-1)}+\frac1{T(B^g-1)}
 \le\frac1{d(B-1)}+\frac1{T(B-1)}.
\]
Across dyadic lengths satisfying $d\le L2^j$, the reciprocal-length
errors sum to at most $2L/[d(B-1)]$.
When $d>2LT$, there is no wrap and $2Lm\le d-1$; hence
$\sum_{i=0}^{d-1}B^i\ge dB^{(d-1)/2}\ge dB^{Lm}$.
Each atom is then at most $1/[d(B-1)]$.
Finally, at most one dyadic length satisfies $LT<d\le2LT$.
For that length the geometric sum and convexity give
\[
 \frac1T\sum_{m=1}^T w_{B,d}(Lm)
 =\frac{B^L}{T(B^L-1)}\frac{B^{LT}-1}{B^d-1}
 \le\frac{LB^L}{d(B^L-1)}
 \le\frac{2L}{d(B-1)}.
\]
For the last inequality, use
$(B^L-1)/(B-1)=\sum_{i=0}^{L-1}B^i\ge B^{L-1}$ and $B\le2$.
Summing the main terms and the two error bounds proves~\eqref{eq:mixed-finite-kernel}.
The release snapshot contains a proof body for the full estimate in
\href{https://github.com/wcook04/plectis-erdos-lean/blob/52f29ad173b04e3bac941b3663f2b9aebe5de0bb/ErdosProblems/Erdos257/PaperCompleteR8/DyadicKernel.lean#L207}{the finite averaging estimate}.
Its complete-cycle and no-wrap scalar ingredients are in
\href{https://github.com/wcook04/plectis-erdos-lean/blob/52f29ad173b04e3bac941b3663f2b9aebe5de0bb/ErdosProblems/Erdos257/PaperCompleteR7/CoverKernel.lean#L130}{the scalar averaging estimates}.
Their formal-checking status is described in Appendix~\ref{app:sources}.
Nonnegative interchange permits summation against any coefficients $c_d$
with $\sum_dc_d/d<\infty$.

\subsection{Positive divisor majorants}
For a finite set $F$, write $f_F(n)=\#\{a\in F:a\mid n\}$.
We bound a fractional power of this count by a nonnegative sum over
divisors of $n$.  The inequality must hold for every positive integer
$n$, not just on average.

\begin{theorem}[a summable divisor-cover criterion]\label{thm:variable-fractional-cover}
For each $j\ge1$, let $F_j\subseteq\Npos$ be finite, let
$0<\alpha_j\le1$, and let $c_{j,d}\ge0$ satisfy
\[
 f_{F_j}(n)^{\alpha_j}\le\sum_{d\mid n}c_{j,d}\quad(n\ge1).
\]
Set $C_j=\sum_{d\ge1}c_{j,d}/d$.  If
\begin{equation*}
 \sum_{j\ge1}\frac{C_j2^{j\alpha_j}}{2^{\alpha_j}-1}<\infty,
 \tag{V}\label{eq:strengthened-cover}
\end{equation*}
then $X_A(b)$ is irrational for every infinite
$A\subseteq\bigcup_jF_j$ and every integer $b\ge2$.
\end{theorem}

Every finite set admits a majorant of the required kind: take
$\alpha_j=1$ and $c_{j,d}=\mathbf1_{F_j}(d)$.  The restriction is that
the costs of the whole sequence satisfy~\eqref{eq:strengthened-cover}.
For a simple admissible family, take $F_j=\{4^j\}$, $\alpha_j=1$, and
$c_{j,4^j}=1$, with all other coefficients zero.  Then $C_j=4^{-j}$ and
the series in~\eqref{eq:strengthened-cover} is $\sum_j2^{-j}$.
The logarithmic obstruction below gives a necessary condition for such a
cover.

\begin{proof}
Write $B_j=2^{\alpha_j}$ and
\[
 U_j(N)=\sum_{r\ge1}2^{-r}f_{F_j}(N+r),\qquad
 V_j(N)=\sum_{d\ge1}c_{j,d}w_{B_j,d}(N).
\]
The choice $B_j=2^{\alpha_j}$ matches the decay after taking a
fractional power.  Subadditivity and the divisor majorant give
\[
 U_j(N)^{\alpha_j}
 \le\sum_{r\ge1}B_j^{-r}f_{F_j}(N+r)^{\alpha_j}
 \le\sum_{d\ge1}c_{j,d}
       \sum_{\substack{r\ge1\\d\mid N+r}}B_j^{-r}
 =V_j(N).
\]
The last identity sums a geometric progression in each residue class.
The geometric-series identity also gives
\[
 \Delta_{2,F_j}(N)=U_j(N)-X_{F_j}(2)\le U_j(N).
\]
Consequently, once the first $J$ finite sets have zero displacement, nonnegativity
bounds the displacement of their union by $\sum_{j>J}U_j(N)$.

Fix $\varepsilon>0$, set $t_j=\varepsilon2^{-j}$, and choose $J$ with
\[
 K_J:=\sum_{j>J}\frac{C_jt_j^{-\alpha_j}}{B_j-1}<\frac14.
\]
This is possible since $\varepsilon^{-\alpha_j}\le\max(1,\varepsilon^{-1})$.
Choose $L$ divisible by every member of the first $J$ finite sets.  Equation~\eqref{eq:mixed-finite-kernel} bounds the finite mean of
$S_J(N):=\sum_{j>J}t_j^{-\alpha_j}V_j(N)$ by
$(1+4L/M)K_J<1/2$ whenever $M\ge4L$.
One sample therefore has $S_J(N)<1$, forcing $U_j(N)<t_j$ for every $j>J$.
Every exponent in the first $J$ finite sets divides $L$, so those terms
have zero displacement.  Consequently,
\[
 0<\Delta_{2,A}(N)\le\sum_{j>J}U_j(N)\le\varepsilon.
\]
The fixed rational lattice excludes rationality at base two.
For $0\le r<d$, the function $(b^r-1)/(b^d-1)$ is nonincreasing
on $b>1$.  For $r>0$, cancel $b-1$ and write it as
$A(b)/(A(b)+C(b))$, where $A(b)=\sum_{i<r}b^i$ and
$C(b)=\sum_{r\le j<d}b^j$.  Its derivative has the sign of
$A'C-AC'=\sum_{i<r\le j<d}(i-j)b^{i+j-1}<0$.
Thus $\Delta_{b,A}(N)\le\Delta_{2,A}(N)$ for every real $b\ge2$;
when $b$ is an integer, the same rational-lattice argument applies.
Any prescribed positive divisor can be included in $L$, so the witnesses
can also be required to be arbitrarily large.
\end{proof}

The same proof permits any positive weights $\eta_j$ with $\sum_j\eta_j=1$:
replace $2^{j\alpha_j}$ in~\eqref{eq:strengthened-cover} by
$\eta_j^{-\alpha_j}$ and take $t_j=\varepsilon\eta_j$.
The earlier condition
$\sum_jC_j2^{j\alpha_j}2^{\alpha_j}/(2^{\alpha_j}-1)^2<\infty$
implies~\eqref{eq:strengthened-cover}, since each old summand is the
new one multiplied by $2^{\alpha_j}/(2^{\alpha_j}-1)>1$.
This proves inclusion of the hypotheses for a given cover, not strict
inclusion of the resulting support classes after optimisation over
all covers.  No separating construction is proved here.

\subsection{What every positive cover must pay}
A necessary lower bound on the cost holds for every choice of cover.  For finite $F$, write
$\mathbb E_F$ for the uniform mean modulo $\operatorname{lcm}(F)$,
with $\operatorname{lcm}(\varnothing)=1$ and $f_{\varnothing}=0$.
More generally allow weights $\eta_j>0$ with $\sum_j\eta_j=1$, and set
\[
 K=\sum_j\frac{C_j\eta_j^{-\alpha_j}}{2^{\alpha_j}-1},\qquad
 \Psi(0)=0,\quad
 \Psi(t)=\inf_{0<\alpha\le1}\frac{t^\alpha}{2^\alpha-1}\quad(t\ge1).
\]
Here $\log^+t=\log\max\{1,t\}$.  For every finite $F$ contained in
$\bigcup_jF_j$,
\begin{equation}
 K\ge\mathbb E_F\Psi(f_F)\ge e\,\mathbb E_F\log^+f_F.
 \label{eq:cover-log-obstruction}
\end{equation}
Indeed, at a point where $f_F(n)=t>0$, some covering set $F_j$ satisfies
$f_{F_j}(n)\ge\eta_jt$.  Otherwise summing contradicts coverage.
The corresponding weighted majorant is at least $\Psi(t)$.
Average first over $1\le n\le X$, using
$\lfloor X/d\rfloor/X\le1/d$, and let $X\to\infty$.
The periodic left side tends to $\mathbb E_F\Psi(f_F)$, proving
the first inequality; the covering moduli need not divide
$\operatorname{lcm}(F)$.  Convexity gives $2^\alpha-1\le\alpha$, and
$e^u/u\ge e$ proves the second.  Those scalar steps have corresponding source bodies in
\href{https://github.com/wcook04/plectis-erdos-lean/blob/52f29ad173b04e3bac941b3663f2b9aebe5de0bb/ErdosProblems/Erdos257/PaperCompleteR7/CoverKernel.lean#L170}{the scalar averaging estimates, lines 170--214},
including the bound $t^\alpha/(2^\alpha-1)\ge e\log t$.
This does not by itself formalise the assembled averaging argument, and
no fresh Lean verification is claimed.

This bound survives optimisation over all covers.  For
$F(q,P)=\{qd:d\mid\prod_{p\in P}p\}$, where $q\ge2$ and no $p\in P$
divides $q$, put $S=\sum_{p\in P}1/p$.  If $S\ge1$, the infimum $K_*$
over finite or countable covers satisfies
\begin{equation}
 \frac{e(S-1)}q\le K_*\bigl(F(q,P)\bigr)\le\frac{eS}q.
 \label{eq:optimal-cube-cost}
\end{equation}
For the lower bound, condition on $q\mid n$ and write
$f_F(n)=2^Z$.  The Chinese remainder theorem gives
$\mathbb EZ=S$.  For $z\ge0$ and $v=\alpha\log2>0$,
\[
 \frac{2^{\alpha z}}{2^\alpha-1}
 \ge\frac{e^{(z-1)v}}v\ge e(z-1)\quad(z>1);
\]
for $0\le z\le1$ the claimed lower bound is nonpositive.
Thus $\Psi(2^z)\ge e(z-1)$.  For the upper bound, use the
one-set cover with $z=1/S$ and $\alpha=\log_2(1+z)$; its exact
positive expansion has cost
\[
 \frac1{qz}\prod_{p\in P}(1+z/p)\le\frac{eS}q.
\]
In particular,
$1-1/S\le qK_*(F(q,P))/(eS)\le1$.
Thus $K_*(F(q,P))\sim eS/q$ as $S\to\infty$, uniformly over the
permitted choices of $q$ and $P$.

The logarithmic lower bound is not a converse: replacing a fractional
majorant by a logarithmic one need not control averages along the
multiples of a prescribed modulus.  A counterexample and the distinct
bound for ordinary initial intervals are proved in the companion,
Section~13, and in~\cite[Theorem~1, Corollary~3 and
Proposition~4]{endpoint2026}.  Those arguments concern finite
functionals, not a comparison of the infinite-support irrationality
classes.

The source
\href{https://github.com/wcook04/plectis-erdos-lean/blob/52f29ad173b04e3bac941b3663f2b9aebe5de0bb/ErdosProblems/Erdos257/PaperCompleteR8/AnalyticSeparationReturn.lean#L16}{constructs a weighted support with no strengthened cover}.
The reverse separation is not established here, so two-way
incomparability is not asserted.  This source statement is distinct
from the finite-functional counterexample just cited.

\subsection{Combining the two support criteria}
Separate small-displacement witnesses need not occur at the same index.
For example, a sequence small only at even indices and one small only at odd
indices need never have a small sum.  The useful feature of (S) is its
uniformity in the moving modulus: the positive-cover argument can use the
exact observation window selected by the weighted proof.
\begin{theorem}[mixed weighted and cover supports]\label{res:mixed-supports}
Let $E,V\subseteq\Npos$.  Suppose $E$ satisfies
\eqref{eq:weighted-return} for a finite nonempty prime set $P$, and
$V\subseteq\bigcup_jF_j$ for finite sets and nonnegative majorants
satisfying the hypotheses of Theorem~\ref{thm:variable-fractional-cover},
with either~\eqref{eq:strengthened-cover} or its positive-weight variant.
Then $X_A(b)$ is irrational for every
infinite $A\subseteq E\cup V$ and every integer $b\ge2$.
\end{theorem}
\begin{proof}
Fix $\varepsilon>0$ and an integer $N_0\ge1$, and put $\rho=\varepsilon/3$.
Use the cover notation $B_j,U_j,V_j$ above, with $\eta_j=2^{-j}$ in
the case~\eqref{eq:strengthened-cover}.  Set $t_j=\rho\eta_j$ and choose
$J$ so that
\[
 K_J=\sum_{j>J}\frac{C_jt_j^{-\alpha_j}}{B_j-1}<\frac1{16}.
\]
Choose a finite $F\subseteq E$ whose complementary weighted mass is
$\kappa<\rho/16$.  Let $L$ be a positive common multiple of $N_0$, all
members of $F$, and all members of the first $J$ finite sets.
For large $H$, choose the same modulus and averaging length as in the base-two weighted proof:
\[
 Q=L\prod_{p\in P}p^{\lfloor\log_pH\rfloor},\qquad
 G=\lfloor H/\max P\rfloor,\qquad M=\lfloor2^{G/2}\rfloor.
\]
The finite estimate~\eqref{eq:weighted-main-bound}, with complementary
weighted mass $\kappa$, remains valid for this $L$: its proof requires
only that every member of $F$ divide $L$.  It also applies to finite or empty $E$, since positivity was
used only after the averaging estimate.  Thus
\[
 \mathscr D_{Q;M,M}(\Delta_{2,E}/\rho)<\frac18
\]
for sufficiently large $H$.  For the same finite distribution, (S) and
nonnegative interchange give
\[
 \mathscr D_{Q;M,M}S_J\le(1+4Q/M)K_J<\frac18,
 \qquad S_J=\sum_{j>J}t_j^{-\alpha_j}V_j,
\]
because $Q/M\to0$.  Hence some sample $N=Qm$ satisfies
$\Delta_{2,E}(N)/\rho+S_J(N)<1$.  At this index the weighted displacement is below $\rho$ and
$U_j(N)<t_j$ for every $j>J$.  The terms with $j\le J$ vanish because
their exponents divide $L$, and hence divide $N$.  Therefore
\[
 \Delta_{2,A}(N)\le\Delta_{2,E}(N)+\Delta_{2,V}(N)<2\rho<\varepsilon,
 \qquad N\ge Q\ge L\ge N_0.
\]
Overlaps between the supports only improve the inequality.  Infinitude of
$A$ makes its displacement positive.  The atom comparison used in the cover
proof transfers arbitrarily small displacements to every integer base;
the fixed rational lattice then excludes rationality.
\end{proof}

The proof does not establish strict containment of either individual
class in the mixed class.

With arbitrary positive cover weights, the class of subsets of such mixed
hosts is closed under finite unions and finite changes.  For weighted
supports use the union of the two finite prime sets: the prime part grows
and $h/(2^h-1)$ decreases.  For two covers, interleave their finite sets with
weights $\eta_j/2$ and $\theta_j/2$; the total cost grows by at most two,
since $2^{\alpha_j}\le2$.  Subsets inherit the same hosts, and finite sets
have finite weighted mass.  This argument uses the positive-weight variant;
it does not silently reindex the dyadic weights in (V).
Countable unions require a tail budget.  Every prime singleton is admitted,
but for the full prime support $\mathcal P$ one has
$\Delta_{2,\mathcal P}(N)>1/3$ for every $N\ge1$: indeed
$\sum_{r\ge1}2^{-r}\omega(N+r)\ge1$, whereas
$X_{\mathcal P}(2)\le\sum_{a\ge2}(2^a-1)^{-1}<2/3$.
Here $\omega(n)$ is the number of distinct prime divisors of $n$.
The reciprocal-summable class is contained in the weighted class, since
$h/(2^h-1)\le1$.  Theorem~\ref{res:reciprocal-support} supplies the direct
proof of that baseline case.


For comparison, Tao--Ter\"av\"ainen prove the full-prime case at base~$2$
\cite[Theorem~1.3, p.~4]{taoteravainen2025}.  The paragraph following it
states extensions to prime support at every integer base and to full
prime-power support at base~$2$, leaving the modifications to the reader.
It does not treat arbitrary infinite thinnings.  The proposed thinning
extension is not a premise of any theorem here.

\section{Finite-support denominator periods}\label{sec:period}

For a finite nonempty $F\subseteq\Npos$, let $D_F$ be the positive reduced
denominator of $X_F(b)=\sum_{n\in F}(b^n-1)^{-1}$.
We use $\operatorname{ord}_1(b)=1$.

\begin{theorem}[the exact denominator period]\label{res:period}
Let $F\subseteq\Npos$ be finite and nonempty, let $b\ge2$ be an integer, and
let $D_F>0$ be the denominator of $X_F(b)$ in lowest terms.  Then $D_F$ is
coprime to $b$, and
\[
  \operatorname{ord}_{D_F}(b)=\lcm\{n:n\in F\}.
\]
If moreover $\lcm(F)\ge2$, then $\lcm(F)<D_F$.  We use
$\operatorname{ord}_1(b)=1$, so the statement includes $F=\{1\}$
at $b=2$.
\end{theorem}

Coprimality is
\mword{Erdos249257/CertificateKernel.lean}{5221}{coprime_base_den_finiteErdosSum}{coprimality of the base and finite-sum denominator}, the order statement is
\mword{Erdos249257/CertificateKernel.lean}{5246}{finite_period_noncollapse_rat_den}{noncollapse for the reduced rational denominator}, and the size bound is
\mword{Erdos249257/CertificateKernel.lean}{5260}{lcm_lt_den_finiteErdosSum}{the lower bound for the finite-sum denominator}.
The three clauses are the CertificateKernel declarations just cited;
\texttt{PaperCompleteR7/Assemblies.lean} is not in this checkout.

Put $L=\operatorname{lcm}(F)$.  Clearing denominators gives
$D_F\mid b^L-1$, so the upper divisibility for the order is immediate.
For the reverse, choose $n\ge2$ maximal under divisibility in $F$ and a
prime $\ell\mid\Phi_n(b)$.  If $e=v_\ell(b^n-1)$, the full prime power
$\ell^e$ has $\operatorname{ord}_{\ell^e}(b)=n$.
Every other selected exponent $m$ has $n\nmid m$, hence
$v_\ell(b^m-1)<e$.  The $n$th summand has uniquely smallest $\ell$-adic
valuation and cannot cancel.  Thus $\ell^e\mid D_F$ and
$n\mid\operatorname{ord}_{D_F}(b)$.  Taking all maximal selected exponents
proves the order statement; the size bound follows from
$\operatorname{ord}_{D_F}(b)\mid\varphi(D_F)<D_F$ when $L\ge2$.
The case $F=\{1\}$ has order one directly.

The same unique-valuation argument permits signs $\pm1$ on the finite
summands.  This is an ordinary deduction here, not a claim of a formal
proof in the unavailable \texttt{SignedFinitePeriodNoncollapse.lean}.
The required cyclotomic prime-power fact is proved in
Section~\ref{app:elementary-details}.
The distinction between primes and prime powers is visible in
\[
 X_{\{2,3\}}(2)=\frac{10}{21},\qquad
 X_{\{2,6\}}(2)=\frac{22}{63}.
\]
In both examples $2$ has multiplicative order six modulo the denominator.
The first combines orders two and three; the second retains the
order-six prime power $9$, although no prime divisor of $63$ has order six.
The boundary $F=\{1\}$ at base two has $D_F=L=1$.
These lower denominator bounds give no upper height control for infinite
partial sums and do not decide an infinite-support value.  By contrast,
Van Assche constructs approximants for the full Lambert series and proves
both nonvanishing and decay of the associated integer linear forms
\cite[Lemma~1 and (34)--(35)]{vanassche2001}.  That approximation mechanism
does not follow from denominator survival alone.

\section{Rational values and scaled tails}\label{sec:forced}
Suppose $X_A(2)=p/v$, where $p\in\Z$ and $v\ge1$ is an integer.
Multiplying the coefficient-series identity by $v2^N$ gives
\[
 z_N:=v\sum_{r\ge1}c_A(N+r)2^{-r}
     =2^Np-v\sum_{n=1}^{N}c_A(n)2^{N-n}\in\Z.
\]
Separating the first term of the tail then gives
$z_{N+1}=2z_N-vc_A(N+1)$.
All the coefficients $c_A(n)$ are divisor counts of the same set $A$.

Two elementary facts delimit what these size estimates can show.
Assume first that $A$ is nonempty and put $a_0=\min A$.  Then every $a_0$ consecutive integers contain a multiple
of $a_0$, so a divisor-incidence zero window has length at most $a_0-1$.
If $A$ is infinite, choose $k$ exponents and a common multiple $L$;
then $c_A(L)\ge k$ and
$\sum_{r\ge1}c_A(L-1+r)2^{-r}\ge k/2$, by its first term.
Both facts hold without rationality.  The arithmetic information lies in
the lattice and in the compatibility of all coefficients with the same
selector.

Precisely, Dirichlet convolution gives
$\mu*c_A=\mathbf1_A$, where $\mu$ is the M\"obius function.
A putative recurrence with integral forcing must therefore satisfy
$(\mu*c_A)(n)\in\{0,1\}$ for all $n$.
The long record gives the integer-recurrence criterion, its telescoping
proof, M\"obius inversion, and the coefficient-series identity in
Theorems~6.105--6.107, under \emph{Bounds for a general coefficient
sequence}.  The example $A=\{2\}$ there has value $1/3$.
The correspondence permits finite supports; infinitude remains a separate
requirement for a counterexample to Problem~\ref{res:problem}.

\section{Limits of the small-displacement criterion}\label{sec:map}
The preceding criteria establish irrationality by making
$\Delta_{2,A}(N)$ positive and arbitrarily small.  The particular
requirement $\inf_{N\ge1}\Delta_{2,A}(N)=0$ fails at full support:
\[
 \Delta_{2,\Npos}(N)
 >(2^N-1)\sum_{a>N}2^{-a}=1-2^{-N}\ge\frac12
 \qquad(N\ge1).
\]
The full-support value is nevertheless
\mword{Erdos249257/CertificateKernel.lean}{8328}{irrational_erdosSum_full_support}{irrational}~\cite{erdos1948}.
Thus a proof covering full support cannot rely solely on this
small-displacement requirement.  This does not exclude other
integer linear forms or other uses of averaging.

The distinction also appears in the squarefree support.  Its series is
irrational at every base $2^j$, $j\ge1$, by Corollary~1.2 and Example~1.1
of Duverney and Tachiya~\cite[p.~4]{duverneytachiya}.  Their corollary
covers the $s$-free products of any pairwise coprime sequence of polynomial
growth, and they present it as support for the conjecture of Erd\H{o}s and
Graham stated here as Problem~\ref{res:problem}.
The zero-window obstruction for one normalised certificate scheme concerns
that scheme's hypotheses.  It gives no obstruction to irrationality of the
value.  The precise normalisation counterexamples are retained in the
long record, together with the complete catalogue of known support families.

A growth-only criterion gives a different class: if
$A=\{c_1<c_2<\cdots\}$ and $\limsup_n c_n/2^n=\infty$, then
$X_A(b)$ is irrational for every integer $b\ge2$
\cite[Theorem~1]{erdos1975}.  The companion, Section~1.7,
proves the translation from denominator growth and gives examples
showing that this class and reciprocal summability are incomparable.
The interval-filling constructions in~\cite{bkkkz2026} and the
freely chosen lacunary denominators of~\cite{vandoornkovac} are not
constructions of subseries of these fixed Mersenne weights.

Prime-incidence arguments require additional care as well: multiplying
all prime exponents by $2$ creates positive correlations between their
divisibility indicators.  The exact identity and covariance calculation
are in the companion, Section~1.5.  Size conditions alone do not
supply that independence.

\section{Unique coding and arithmetic membership}\label{sec:geometry}
Let $w_n=(2^n-1)^{-1}$ and
\[
 \Ach=\left\{\sum_{n\ge1}\varepsilon_nw_n:
                   \varepsilon_n\in\{0,1\}\right\},\qquad
 R_N=\sum_{n>N}w_n.
\]
The estimate
\[
 2^{-N}<R_N\le2^{-N}+\frac23\,4^{-N}<w_N\qquad(N\ge1)
\]
implies uniqueness of the selector.  Indeed, at the first differing
exponent $n$, the difference $w_n$ exceeds everything later terms can
cancel.  The middle inequality follows by writing
$w_n=2^{-n}+4^{-n}/(1-2^{-n})$ and using $n\ge N+1\ge2$.
Since every weight exceeds its tail, Hornich's theorem~\cite{hornich1941}, as proved by
Nitecki~\cite[Theorem~4(1), p.~9]{nitecki2013}, shows that the coding image is
a Cantor set of measure $\lim_N2^NR_N$; the estimate above gives
$2^NR_N\to1$, so $\lambda(\Ach)=1$.
Concretely, fixing the first $N$ digits gives $2^N$ disjoint closed
intervals of length $R_N$.  These interval unions decrease to $\Ach$,
and continuity of Lebesgue measure from above gives the same limit.
Kova\v{c}--Tao record this strict-tail inequality and the Cantor conclusion in the fixed-base setting~\cite[Remark~4.1]{kovactao}.
For a specified rational target, the useful consequence here is
uniqueness of its possible selector.  Restricted-set dimension and
measure calculations are given in the companion, Section~1.8.

\section{Greedy membership and an integer recurrence}\label{sec:actual-repairs}
For $x\ge0$, the greedy rule starts with $r_0=x$ and, at rank $n\ge1$,
selects $n$ when $r_{n-1}\ge(2^n-1)^{-1}$, subtracting that weight if
selected.  Let $A_x$ be the resulting support and
$c_x(n)=\#\{a\in A_x:a\mid n\}$.  Set
\[
 P_0=0,\qquad P_{N+1}=2P_N+c_x(N+1),\qquad
 Q_N=\lfloor2^Nx\rfloor-P_N.
\]
Thus $P_N=\sum_{j=1}^N2^{N-j}c_x(j)$ is the integer truncation of the
Lambert expansion~\eqref{eq:incidence}.  It is at most
$2^NX_{A_x}(2)\le2^Nx$, so $Q_N\ge0$.  These integer remainders obey
\begin{equation}
 Q_{N+1}=2Q_N+\beta_N-c_x(N+1),\qquad
 \beta_N=\lfloor2^{N+1}x\rfloor-2\lfloor2^Nx\rfloor\in\{0,1\}.
 \label{eq:actual-repair-recurrence}
\end{equation}

\begin{theorem}[membership and nonincreasing integer remainders]\label{res:general-repair}
For every real $x\ge0$, the following are equivalent:
\[
 \begin{gathered}
 x\in\Ach;\\
 \forall K\ge0\ \exists N\ge K:\quad Q_{N+1}\le Q_N;\\
 \forall K\ge0\ \exists N\in[K,K+2\lfloor\sqrt K\rfloor+12):
 \quad Q_{N+1}\le Q_N.
 \end{gathered}
\]
\end{theorem}

The following is an ordinary proof.  The supplied solution wrapper
\nolinkurl{Solutions/PalomarCorpus/E257/GeneralRepairCriterion.lean}
contains the two named equivalences, with explicit bridges to the challenge
definitions and the same square-root window.  Its provenance is the
separate release \texttt{52f29ad173b0}.  This static comparison is not a
new replay; a challenge declaration alone is not cited as a proof.

\begin{proof}
Strict domination of each Mersenne weight over its tail implies that a
represented target is recovered by the greedy rule.  For such a target,
\[
 0\le Q_N\le\sum_{r\ge1}c_x(N+r)2^{-r}
 \le2\sqrt N+4,
\]
since $c_x(n)\le\tau(n)\le2\sqrt n$ and
$\sqrt{N+r}\le\sqrt N+\sqrt r\le\sqrt N+r$.
Put $s=\lfloor\sqrt K\rfloor$ and $T=2s+12$.
Strict increase at all $T$ steps would imply $Q_{K+T}\ge T$.
But $K+T<(s+4)^2$ gives $Q_{K+T}<2s+12=T$.
This proves the window condition and hence cofinal nonincreases.

Conversely, if $x\notin\Ach$, then
$\delta=x-X_{A_x}(2)>0$.  The Lambert prefix is at most
$2^NX_{A_x}(2)$, so $Q_N\ge2^N\delta-1$.
Eventually this exceeds $c_x(N+1)\le N+1$.
Equation~\eqref{eq:actual-repair-recurrence} then makes $Q_N$ strictly
increasing, contradicting cofinal nonincreases.
\end{proof}

Theorem~\ref{res:general-repair} characterises membership, but it does not
locate the required nonincreases for a specified rational target.  For $x=0$, the selector is empty and $Q_N=0$ at every
rank.  A finite subseries sum is likewise represented and satisfies the
criterion.  In contrast, $x=3/4$ lies strictly between $R_1<2/3$ and
$w_1=1$, so it is not represented and its integer remainders eventually
increase strictly.  For $1/2$ and $1/21$, the theorem does not establish
which alternative occurs.

The square-root estimate is uniform in the target.  There is also a
uniform subpower refinement: for each $0<\varepsilon<1$, a constant
$C_\varepsilon$ gives the window length
$\lceil C_\varepsilon(K+1)^\varepsilon\rceil$.
Indeed, $\tau(n)\le A_\varepsilon n^\varepsilon$ gives
$Q_N\le D_\varepsilon(N+1)^\varepsilon$ for represented targets,
where $D_\varepsilon=A_\varepsilon\sum_{r\ge1}r^\varepsilon2^{-r}$.
Choose $C_\varepsilon>D_\varepsilon(C_\varepsilon+3)^\varepsilon$,
which is possible because $\varepsilon<1$.
For $T=\lceil C_\varepsilon(K+1)^\varepsilon\rceil$, strict increase
at all $T$ steps would give
\[
 T\le Q_{K+T}\le D_\varepsilon(K+T+1)^\varepsilon
 \le D_\varepsilon(C_\varepsilon+3)^\varepsilon(K+1)^\varepsilon<T,
\]
a contradiction.  The converse still follows from exponential growth
when $x$ is not represented.  This ordinary argument strengthens the
window bound, not the occurrence claim for a specified target; it is
separate from the square-root statement in the cited formal source.

At $x=1/2$, the digits $\beta_N$ vanish for $N\ge1$; at $x=1/21$ they
are six-periodic.  Thus the unresolved arithmetic input is
\begin{equation}
 \boxed{\quad
 \forall K\ \exists N\ge K:\qquad
 c_x(N+1)\ge Q_N+\beta_N,
 \qquad x\in\{1/2,1/21\}.
 \quad}\label{eq:actual-selector-obligation}
\end{equation}
Both $Q_N$ and $c_x$ must arise from the same greedy selector.
The long record retains exact counterexamples to fixed-multiplier repair
schedules and the finite phase masks used to test them.
Neither target is decided here.

\section{Further questions}\label{sec:open}
The support theorems give sufficient conditions for irrationality;
the tests below instead concern whether a specified rational is a
subseries sum.  Their hypotheses must be verified for the actual
selector or for finite approximants at arbitrarily large depths.
An equivalent test does not itself establish membership.

\paragraph{Integer quotients for $1/21$.}
For $M,d\ge1$, set
\[
 q_M(d)=\left\lfloor\frac{2^M}{2^d-1}\right\rfloor,
 \qquad T_M=\left\lfloor\frac{2^M}{21}\right\rfloor.
\]
Starting with the integer remainder $T_{2R}$, consider $d=2,\ldots,R$
in order and subtract $q_{2R}(d)$ whenever it does not exceed the current
remainder.  Let $D_R$ be the set of selected exponents and $s_R$ the final
remainder.  Thus
\[
 s_R=T_{2R}-\sum_{d\in D_R}q_{2R}(d)\ge0.
\]
A second comparison uses the same integer rule through $d=2R$, rather
than stopping at $R$.
For example, $R=6$ gives $T_{12}=195$, $D_6=\{5\}$ and
$s_6=195-\lfloor4096/31\rfloor=63\le64$.  This is one row satisfying
the bound below, not evidence that such rows occur at unbounded depths.


Write $r_n=r_n(1/21)$ for the real greedy remainder and $A_{1/21}$ for its
support.  The condition $\mathcal F_{21}$ means that there exist integers
$n,R_0,K_0\ge0$ with all of the following properties:
\begin{enumerate}
\item $r_n>R_n$, the set of positive exponents outside $A_{1/21}$ is
finite, and every exponent greater than $n$ belongs to $A_{1/21}$;
\item for every $R\ge R_0$, the integer rule with target $T_{2R}$ and
weights $q_{2R}(2),\ldots,q_{2R}(2R)$ selects exactly the same exponents
as the real greedy rule on $\{2,\ldots,2R\}$;
\item for every $K\ge K_0$, the interval $(K,2K]$ contains an exponent
of $A_{1/21}$.
\end{enumerate}
Some of these clauses follow from others once the eventual all-selected
behaviour is known.  They are written out to match the precise branch
used in the formal result.  In particular, $\mathcal F_{21}$ is not a
hypothesis about an arbitrary recurrence with similar coefficients.

\begin{theorem}[integer-quotient tests for $1/21$]
\label{res:one-over-twenty-one-frontier}
The following statements hold.
\begin{enumerate}
\item $1/21\in\Ach$ if and only if $\mathcal F_{21}$ does not hold.
\item If there is an unbounded sequence of ranks $R$ with $s_R\le 2^R$,
  then $1/21\in\Ach$.
\item On $\mathcal F_{21}$, eventually $s_R>2^R$, the boundary rank $R+1$
  belongs to $D_{R+1}$, and, for all sufficiently large $R$,
\[
 \begin{aligned}
 D_{R+1}&=D_R\cup\{R+1\},\\
 s_{R+1}&=4s_R+
 \left\lfloor\frac{4(2^{2R}\bmod21)}{21}\right\rfloor
 -2c_{D_R}(2R+1)-c_{D_R}(2R+2)\\
 &\hspace{3em}{}-(2^{R+1}+1).
 \end{aligned}
\]
Here $c_{D_R}(m)=\#\{d\in D_R:d\mid m\}$.
\end{enumerate}
\end{theorem}

\noindent The equivalence is
\mword{Erdos249257/TwentyOneQuotientGreedy.lean}{3507}{one_div_twenty_one_mem_iff_not_fatalAlignedBranch}{the equivalence with failure of $\mathcal F_{21}$};
the closed-row compactness step is
\mword{Erdos249257/TwentyOneQuotientGreedy.lean}{5554}{twentyOneCofinalEvenQuotientGreedyDecay_of_closedRows}{vanishing scaled error from unbounded closed rows};
and the eventual affine regime is
\mword{Erdos249257/TwentyOneQuotientGreedy.lean}{5658}{twentyOneFatalAlignedBranch_eventually_affine_supercapacity}{the eventual affine recurrence on $\mathcal F_{21}$}.
The finite uniqueness statement is also explicit: if
$D\subseteq\{2,\ldots,R\}$ and an integer $s$ satisfy
$\sum_{d\in D}q_{2R}(d)+s=T_{2R}$ with $0\le s\le2^R$, then
$D=D_R$ and $s=s_R$.  This is the denominator-specific separation theorem
(\mword{Erdos249257/TwentyOneQuotientGreedy.lean}{231}{twentyOneClosedRow_forces_quotientGreedy}{uniqueness of a finite representation with the stated remainder bound}).
The cited finite crossing lemmas give additional consequences under their
alignment hypotheses: an earlier finite prefix cannot occur, and a real
greedy exponent must be skipped
(\mword{Erdos249257/TwentyOneQuotientGreedy.lean}{5179}{twentyOneAlignedSaturatedCrossing_forces_canonical_ancestor_hole}{the missing-prefix consequence of an aligned crossing},
\mword{Erdos249257/TwentyOneQuotientGreedy.lean}{5223}{twentyOneAlignedSaturatedCrossing_forces_scaled_greedy_skip}{the real greedy skip forced by that crossing}).

\paragraph{Approximation to $1/2$ by finite supports.}
For $A\subseteq\{2,\ldots,M\}$ and $1\le m\le M$, define the integer
\[
 K_A(m)=2^{m-1}-\sum_{j=2}^{m}2^{m-j}c_A(j).
\]
It measures the error in the truncated divisor-coefficient sum, after
multiplication by $2^m$.  In particular,
$K_A(1)=1$ and $K_A(m+1)=2K_A(m)-c_A(m+1)$.

\begin{theorem}[finite approximations with vanishing scaled error]
\label{res:terminalhalf}
Suppose there are integers $M_j\ge1$ tending to infinity and sets
$A_j\subseteq\{2,\ldots,M_j\}$ such that
\[
 \frac{|K_{A_j}(M_j)|}{2^{M_j}}\longrightarrow0.
\]
Then $X_A(2)=1/2$ for some infinite set $A\subseteq\Npos$.
\end{theorem}

No agreement between different $A_j$ is assumed, and no bound is imposed
on their earlier carries.  Producing such finite approximants would
refute Problem~\ref{res:problem}; their existence is not established
here.  The implication follows from the estimate
\[
 \left|X_{A_j}(2)-\frac12\right|
 \le \frac{|K_{A_j}(M_j)|+2\sqrt{M_j}+4}{2^{M_j}}.
\]
It follows by separating the Lambert-series coefficient tail and using
$c_{A_j}(n)\le\tau(n)\le2\sqrt n$.  The finite sums therefore tend to
$1/2$.  The achievement set is closed, and no finite support has value
$1/2$, so its representing support is infinite.  A square-root bound on
the terminal carry would suffice, but the theorem permits any error
$o(2^{M_j})$.  For a finite $D\subseteq\{2,\ldots,M\}$, the
floor-quotient identity gives
\[
 K_D(M)=2^{M-1}-\sum_{d\in D}
                \left\lfloor\frac{2^M}{2^d-1}\right\rfloor.
\]
Thus an exact quotient row has $K_D(M)=1$, whereas the terminal
criterion permits a larger error of either sign.  These are different
conditions at a fixed depth, even though their cofinal existence
conditions both characterise half-membership.

Conversely, if $X_A(2)=1/2$, then $1\notin A$, and the prefixes
$A_M=A\cap\{2,\ldots,M\}$ satisfy
\[
 K_{A_M}(M)=\sum_{r\ge1}c_A(M+r)2^{-r}
       \le 2\sqrt M+4.
\]
The equality uses the coefficients of the full support $A$ on the
right: truncation does not change the coefficients through $M$.
Hence the existence hypothesis of Theorem~\ref{res:terminalhalf}
is equivalent to half-membership.  Allowing incompatible finite
supports removes a construction requirement, not the difficulty of
the existence problem.  Compactness proves the implication; it does
not produce the approximating supports.

\paragraph{Finite families with a shared prefix.}
Put $B(m)=2\lfloor\sqrt m\rfloor+4$.  For $0\le K\le M$, consider a
family $A_1,\ldots,A_{B(M)}\subseteq\{2,\ldots,M\}$ with
\[
 1\le K_{A_k}(m)\le B(m)\quad(1\le m\le M),
 \qquad K_{A_k}(M)=k.
\]
Require that all the $A_k$ agree on $\{1,\ldots,K\}$ and that some
integer $E\ge B(M)$ satisfies
\[
 \sum_{j=K+1}^{M}1_{A_k}(j)\,2^{M-j}+k=E
 \qquad(1\le k\le B(M)).
\]
Thus the binary suffixes have consecutive values $E-k$, while the
terminal carry ranges through the entire permitted interval.  This is
substantially more data than one finite approximation to $1/2$.

\begin{theorem}[unbounded shared-prefix families imply a half-support]
\label{res:cylinderhalf}
Suppose that for every $N$ there are $M,K$ with
$\max\{N,1\}\le M$, $0\le K\le M$, and a family satisfying all the
conditions in the preceding paragraph.  Then $X_A(2)=1/2$ for some
infinite set $A\subseteq\Npos$.
\end{theorem}

Choosing one member from each family gives the terminal bound in
Theorem~\ref{res:terminalhalf}, proving the conclusion.  The shared-prefix
conditions are needed for the proposed construction of the families,
not for this final compactness step.  Such a construction at unbounded
depths would refute Problem~\#257; none is proved here.

\begin{problem}[membership of 1/21 in the Mersenne achievement set]
\label{prob:one-over-twenty-one-membership}
For the greedy remainder of $1/21$, prove that arbitrarily large $N$
satisfy
\[
 2^{N+1}r_N(1/21)<\frac{2^{N+1}}{2^{N+1}-1},
 \quad\text{equivalently}\quad
 r_N(1/21)<\frac1{2^{N+1}-1}.
\]
Equivalently, exclude the condition $\mathcal F_{21}$ defined above.
One sufficient route is to rule out the eventual recurrence in
Theorem~\ref{res:one-over-twenty-one-frontier}(3) together with $s_R>2^R$;
another is to prove $s_R\le2^R$ at arbitrarily large ranks.  Neither
sufficient route is claimed to be necessary by itself.
\end{problem}

\begin{problem}[bounded returns of the scaled remainder]\label{prob:scaled-return}
For each of the targets $x=1/2$ and $x=1/21$, does the greedy remainder
$r_N(x)$ return to one bounded interval after scaling by $2^N$?
\[
 \exists B<\infty\ \forall K\ \exists N\ge K:
 \qquad 2^N r_N(x)\le B.
\]
\end{problem}

\begin{problem}[arithmetic tests for the greedy sequence]\label{prob:actual-invariant}
Can a finite-memory, $2$-adic or discrepancy argument prove $s_R\le2^R$
at arbitrarily large ranks, or rule out the eventual recurrence in
Theorem~\ref{res:one-over-twenty-one-frontier}(3)?  Such an argument must use
the divisor counts of the actual greedy support.  Can a bounded window of
$R\bmod6$, residues of $s_R$, endpoint divisor counts and the finite set of
eventual skips force a decrease or a contradiction?  Alternatively, can one
show that these bounded-memory data cannot distinguish the actual sequence
from sequences that remain above $2^R$ but need not come from a support?
\end{problem}

\begin{problem}[final-skip Diophantine exclusion]\label{prob:fatal-interval}
Let $E=\sum_{n\ge1}(2^n-1)^{-1}$.  If $1/21\notin\Ach$, let $M$ be the last
skipped exponent, $S_M$ its finite skipped prefix, and
\[
 a_M=\frac1{21}+\sum_{d\in S_M}\frac1{2^d-1}.
\]
Write $\operatorname{gap}_M=(2^M-1)^{-1}-R_M>0$.
Can the arithmetic restrictions on this last skipped exponent be used to
prove $|E-a_M|\ge\operatorname{gap}_M$, contradicting
\[
 0<a_M-E<\operatorname{gap}_M?
\]
\end{problem}


\section{Elementary details used in the proofs}\label{app:elementary-details}
\paragraph{Cyclotomic prime-power fact.}
Let $b\ge2$, $n\ge2$, $\ell\mid\Phi_n(b)$ be prime and
$e=v_\ell(b^n-1)$.  Then $\operatorname{ord}_{\ell^e}(b)=n$.
Here is a proof including the exceptional $2$-adic behaviour.
For odd $\ell$, put $d=\operatorname{ord}_{\ell}(b)$ and
$s=v_\ell(b^d-1)$.  The elementary lifting identity
$v_\ell(b^{dt}-1)=s+v_\ell(t)$ follows by factoring a geometric sum
when $\ell\nmid t$ and by a binomial expansion for a factor $\ell$.
In the factorisation $b^n-1=\prod_{r\mid n}\Phi_r(b)$, subtracting
these valuations over proper divisors gives
\[
 v_\ell(\Phi_n(b))=
 \begin{cases}s&n=d,\\1&n=d\ell^j,\ j\ge1,\\0&\text{otherwise.}\end{cases}
\]
If $n=d$, no smaller positive exponent gives $b^m\equiv1\pmod\ell$.
If $n=d\ell^j$, the full valuation is $e=s+j$; the same lifting
identity shows that the least exponent giving valuation $e$ is
$d\ell^j=n$.
For $\ell=2$, $b$ is odd.  Factoring $b^{2t}-1$ gives
\[
 v_2(b^{2^j}-1)=v_2(b-1)+v_2(b+1)+j-1\qquad(j\ge1).
\]
The only possible indices $n\ge2$ are $n=2^j$:
$v_2(\Phi_2(b))=v_2(b+1)$ and
$v_2(\Phi_{2^j}(b))=1$ for $j\ge2$; all other indices have
valuation zero.  The displayed identity again makes $2^j$ the least
exponent with the full valuation.  This proves the fact.
Finally $\Phi_n(b)>1$: each primitive $n$th root $\zeta$ satisfies
$|b-\zeta|>1$, and their product is $\Phi_n(b)$.

\paragraph{Exact scalar cover cost.}
For $t\ge1$, put $z=\log_2t$.  The substitution $y=2^\alpha$
reduces the infimum defining $\Psi(t)$ to $y^z/(y-1)$ on $1<y\le2$.
Its derivative has the sign of $(z-1)y-z$.  Hence
\[
 \Psi(t)=
 \begin{cases}t,&1\le t\le4,\\
 z^z/(z-1)^{z-1},&t>4.
 \end{cases}
\]
The endpoint and the open lower bound $\alpha>0$ are both included in this
calculation.  This scalar identity does not supply a converse to the
positive-cover criterion.

\appendix
\section{Guide to the formal sources}\label{app:sources}

The public proof closure at \texttt{065e0952d7a3} was replayed under
Lean~4.29.1.  Besides the reciprocal and weighted declarations cited above,
the replay checked the exact strengthened positive-cover conclusion in
\href{https://github.com/wcook04/plectis-erdos/blob/065e0952d7a3a0c9f87275323d50aaedb78f481f/lean/ErdosProblems/Erdos257/PaperCompleteR8/PositiveCoverReturn.lean#L241}{\texttt{strengthenedPositiveCoverClaim}}
and the common-witness conclusion in
\href{https://github.com/wcook04/plectis-erdos/blob/065e0952d7a3a0c9f87275323d50aaedb78f481f/lean/ErdosProblems/Erdos257/PaperCompleteR8/WeightedReturn.lean#L105}{\texttt{mixedSupportClaim}}.
The axiom audit for these endpoints reports only \texttt{propext},
\texttt{Classical.choice}, and \texttt{Quot.sound}.  The ordinary proof for
$A_\star$ remains separate from the existential host wrapper, which does not
verify that named example.

Erd\H{o}s's five-page paper was checked directly: p.~222 states the
reciprocal-summable extension without proof, and p.~226 describes the
fractional-part approach.  The original Luca--Tachiya and Hornich articles
were not independently retrieved.  The periodic theorem
\cite{lucatachiya2014periodic} was checked in Luca and Tachiya's own
account~\cite[Theorem~A and Example~2, pp.~139--140]{lucatachiya2017},
and the strict-tail result in Nitecki's exposition.  The \emph{Formal Conjectures}
file~\cite{formalconjectures257} is statement-level prior art, not a
proof dependency.

The index retains the original names and revisions: standard macros use
\texttt{\commitshort}, and fourteen paper-local coordinates across the two
papers use \texttt{f36a98bf3d3e}.  The linked hypotheses and conclusions,
not the abbreviated index labels, give the precise statements.

\paragraph{Finite sums and rational denominators.}
\mword{Erdos249257/CertificateKernel.lean}{5091}{finite_period_noncollapse}{noncollapse of the finite period};
\mword{Erdos249257/CertificateKernel.lean}{5246}{finite_period_noncollapse_rat_den}{noncollapse for the reduced rational denominator};
\mword{Erdos249257/CertificateKernel.lean}{5221}{coprime_base_den_finiteErdosSum}{coprimality of the base and finite-sum denominator};
\mword{Erdos249257/CertificateKernel.lean}{5260}{lcm_lt_den_finiteErdosSum}{the lower bound for the finite-sum denominator}.

\paragraph{Rationality and integer carries.}
\mword{Erdos249257/GenericTailOrbitRigidity.lean}{426}{binaryCoeffSeries_rational_iff_exists_temperedBinaryOrbit}{rationality and integer recurrences with vanishing scaled error};
\mword{Erdos249257/RationalSupportCarrySkeleton.lean}{2383}{exists_unbounded_shifted_odd_tail_nat_state_of_support_fraction}{unboundedness of the shifted odd-tail recurrence};
\mword{Erdos249257/SublogDivisorCoverage.lean}{392}{supportCoeffZeroWindow_length_le_eps_logb_add}{lengths of intervals with zero divisor count};
\mword{Erdos249257/RationalSupportCarrySkeleton.lean}{1480}{one_div_oddOrder_le_reciprocalMass_of_support_fraction}{a lower bound for reciprocal mass};
\mword{Erdos249257/RationalSupportCarrySkeleton.lean}{2210}{dyadic_support_fraction_reciprocalMass_diverges_or_gt_one}{the reciprocal-mass alternative for dyadic values};
\mword{Erdos249257/BooleanMobiusCarry.lean}{949}{exists_normalized_support_fraction_iff_exists_booleanMobiusCarry}{rational support values and M\"obius-inverted carries}.

\paragraph{Classical support theorems and examples.}
The formal
\mword{Erdos249257/CertificateKernel.lean}{10776}{irrational_erdosSupportSeries_pairwise_coprime}{pairwise-coprime theorem}
retains the coprimality hypothesis, whereas
\mword{Erdos249257/CertificateKernel.lean}{8328}{irrational_erdosSum_full_support}{full-support irrationality}
is a separate result.  The other classical-support declarations are:
\mword{Erdos249257/CertificateKernel.lean}{9045}{irrational_erdosSupportSeries_univ}{irrationality for full support};
\mword{Erdos249257/CertificateKernel.lean}{9103}{irrational_erdosSupportSeries_multiples}{irrationality for multiples of an integer};
\mword{Erdos249257/CertificateKernel.lean}{12811}{irrational_ratWeightSeries_eventuallyPeriodic}{irrationality for the stated eventually periodic weights};
\mword{Erdos249257/CertificateKernel.lean}{11672}{irrational_erdosSupportSeries_residueClass}{residue class};
\mword{Erdos249257/CertificateKernel.lean}{11686}{irrational_erdosSupportSeries_odd}{odd};
\mword{Erdos249257/CertificateKernel.lean}{6035}{irrational_erdosSum_factorial_support}{factorial support};
\mword{Erdos249257/CertificateKernel.lean}{6059}{irrational_erdosSum_two_pow_support}{powers-of-two support};
\mword{Erdos249257/CertificateKernel.lean}{6082}{erdos257_family_factorial_instance}{factorial-support instance};
\mword{Erdos249257/CertificateKernel.lean}{6090}{erdos257_family_two_pow_instance}{powers-of-two-support instance};
\mword{Erdos249257/CertificateKernel.lean}{10776}{irrational_erdosSupportSeries_pairwise_coprime}{pairwise-coprime support};
\mword{Erdos249257/CertificateKernel.lean}{8328}{irrational_erdosSum_full_support}{full-support irrationality};
\mword{Erdos249257/CertificateKernel.lean}{6272}{lcm_gap_hypothesis_fails_full_support}{failure of the least-common-multiple gap condition for full support};
\mword{Erdos249257/CertificateKernel.lean}{14175}{irrational_or_bpow_mul_eq_intCast_intWeightedErdosSeries_periodic}{the rationality alternative for signed periodic weights}.

\paragraph{Squarefree divisor counts and certificate restrictions.}
\lword{Erdos257/SquarefreeSupportIncidence.lean}{94}{card_squarefreeDivisors}{the number $2^{\omega(n)}$ of squarefree divisors};
\lword{Erdos257/SquarefreeSupportIncidence.lean}{111}{squarefreeIncidence_eq}{the squarefree divisor-count identity};
\lword{Erdos257/SquarefreeSupportIncidence.lean}{140}{odd_squarefreeIncidence}{parity of the squarefree divisor count};
\lword{Erdos257/SquarefreeSupportIncidence.lean}{71}{squarefreeDivisors_eq_image}{squarefree divisors as products of distinct primes};
\lword{Erdos257/SquarefreeSupportIncidence.lean}{274}{not_exists_carry_certificates_squarefreeSupport}{failure of the stated carry certificates for squarefree support};
\lword{Erdos257/SquarefreeSupportIncidence.lean}{292}{not_exists_digitwise_certificates_squarefreeSupport}{failure of the stated digitwise certificates for squarefree support};
\lword{Erdos257/SquarefreeSupportIncidence.lean}{314}{supportCoeff_squarefreeShiftedSupport}{$2^{\omega(n)}$ incidence};
\lword{Erdos257/SquarefreeSupportIncidence.lean}{335}{irrational_erdosSupportSeries_squarefreeSupport_iff_shifted}{equivalence after shifting the squarefree divisor count};
\lword{Erdos257/SquarefreeSupportIncidence.lean}{441}{exists_omega_ge_block}{blocks on which $\omega(n)$ is large};
\lword{Erdos257/SquarefreeSupportIncidence.lean}{485}{exists_digitwise_block_squarefreeShiftedSupport}{digitwise blocks for the shifted divisor count}.

\paragraph{Deleting finitely many terms.}
\mword{Erdos249257/CertificateKernel.lean}{9467}{irrational_erdosSupportSeries_of_tail}{passing irrationality from a tail to the series};
\mword{Erdos249257/CertificateKernel.lean}{9476}{irrational_erdosSupportSeries_tail_of_irrational}{passing irrationality from the series to a tail}.

\paragraph{The full achievement set.}
\mword{Erdos249257/GreedyAchievementSet.lean}{656}{isCompact_mersenneAchievementSet}{compact};
\mword{Erdos249257/GreedyAchievementSet.lean}{1656}{perfect_mersenneAchievementSet}{perfect};
\mword{Erdos249257/GreedyAchievementSet.lean}{1672}{isTotallyDisconnected_mersenneAchievementSet}{totally disconnected};
\mword{Erdos249257/GreedyAchievementSet.lean}{1681}{isNowhereDense_mersenneAchievementSet}{nowhere dense};
\mword{Erdos249257/GreedyAchievementSet.lean}{996}{volume_mersenneAchievementSet}{Lebesgue measure one};
\mword{Erdos249257/GreedyAchievementSet.lean}{1458}{mem_mersenneAchievementSet_iff_greedy_survival}{membership and greedy tail bounds};
\mword{Erdos249257/GreedyAchievementSet.lean}{1784}{three_fourths_not_mem_mersenneAchievementSet}{non-membership of $3/4$}.

\paragraph{Restricted achievement sets.}
\lword{Erdos257/MersenneSubseriesRigidity.lean}{20}{selectedMersenneTail}{the tail of a restricted subseries};
\lword{Erdos257/MersenneSubseriesRigidity.lean}{24}{summable_selectedMersenneTail}{summability of the restricted tail};
\lword{Erdos257/MersenneSubseriesRigidity.lean}{30}{selectedMersenneTail_lt_weight}{strict domination of the restricted tail by the preceding weight};
\lword{Erdos257/MersenneSubseriesRigidity.lean}{54}{supportedMersenneDigitValue_injective}{injectivity of the restricted digit map};
\lword{Erdos257/MersenneSubseriesRigidity.lean}{45}{SupportedMersenneDigits}{digit strings vanishing off $J$};
\lword{Erdos257/MersenneSubseriesRigidity.lean}{49}{supportedMersenneDigitValue}{restricted digit map};
\lword{Erdos257/MersenneSubseriesRigidity.lean}{63}{supportedMersenneDigitSet}{supported digit set};
\lword{Erdos257/MersenneSubseriesRigidity.lean}{66}{isClosed_supportedMersenneDigitSet}{closed};
\lword{Erdos257/MersenneSubseriesRigidity.lean}{76}{supportedMersenneAchievementSet}{restricted achievement set};
\lword{Erdos257/MersenneSubseriesRigidity.lean}{79}{supportedMersenneAchievementSet_eq_image}{the image theorem};
\lword{Erdos257/MersenneSubseriesRigidity.lean}{90}{isCompact_supportedMersenneAchievementSet}{compact};
\lword{Erdos257/MersenneSubseriesRigidity.lean}{96}{isClosed_supportedMersenneAchievementSet}{closed};
\lword{Erdos257/MersenneSubseriesRigidity.lean}{103}{supportedMersenneAchievementSet_subset}{support restriction};
\lword{Erdos257/MersenneSubseriesRigidity.lean}{112}{isNowhereDense_supportedMersenneAchievementSet}{nowhere dense};
\lword{Erdos257/MersenneSubseriesRigidity.lean}{120}{preperfect_supportedMersenneDigitSet}{absence of isolated restricted digit strings for infinite support};
\lword{Erdos257/MersenneSubseriesRigidity.lean}{150}{preperfect_supportedMersenneAchievementSet}{absence of isolated represented values for infinite support};
\lword{Erdos257/MersenneSubseriesRigidity.lean}{167}{perfect_supportedMersenneAchievementSet}{perfect};
\lword{Erdos257/MersenneSubseriesRigidity.lean}{175}{summable_mersenneDigitTerm}{controls digit terms};
\lword{Erdos257/MersenneSubseriesRigidity.lean}{182}{positiveMersenneDigitValue_update}{the update formula};
\lword{Erdos257/MersenneSubseriesRigidity.lean}{203}{supportedMersenneAchievementSet_insert}{the union formula after adding one exponent};
\lword{Erdos257/MersenneSubseriesRigidity.lean}{262}{disjoint_supportedMersenneAchievementSet_translate}{disjointness of the two translated copies};
\lword{Erdos257/MersenneSubseriesRigidity.lean}{286}{volume_supportedMersenneAchievementSet_insert}{doubles the volume};
\lword{Erdos257/MersenneSubseriesRigidity.lean}{300}{supportedMersenneAchievementSet_univ}{recovery of the full achievement set};
\lword{Erdos257/MersenneSubseriesRigidity.lean}{313}{volume_supportedMersenneAchievementSet_finset_compl_scaled}{the scaled volume identity after deleting finitely many exponents};
\lword{Erdos257/MersenneSubseriesRigidity.lean}{349}{volume_supportedMersenneAchievementSet_finset_compl}{gives $2^{-|F|}$};
\lword{Erdos257/MersenneSubseriesRigidity.lean}{358}{supportedMersenneAchievementSet_mono}{monotonicity under inclusion of supports};
\lword{Erdos257/MersenneSubseriesRigidity.lean}{368}{volume_supportedMersenneAchievementSet_eq_zero_of_compl_infinite}{measure zero};
\lword{Erdos257/MersenneSubseriesRigidity.lean}{397}{volume_supportedMersenneAchievementSet_dichotomy}{the formal dichotomy}.

\paragraph{Half-membership and finite approximations.}
\mword{Erdos249257/GreedyAchievementSet.lean}{2583}{half_mem_mersenneAchievementSet_iff_greedySkippedSupport_infinite}{the greedy form};
\mword{Erdos249257/HalfCylinderHalfMembershipClassification.lean}{126}{half_mem_mersenneAchievementSet_iff_unboundedTerminalFalse}{the terminal-bit form};
\mword{Erdos249257/HalfCylinderHalfMembershipClassification.lean}{213}{half_mem_mersenneAchievementSet_iff_exists_unboundedSkippedRanksAlong}{the skipped-rank form};
\mword{Erdos249257/HalfCylinderFatalGapRightTail.lean}{781}{seamGreedyEventuallyRight_iff_existsFatalHalfGap}{the fatal-gap equivalence};
\mword{Erdos249257/HalfCylinderFatalGapRightTail.lean}{787}{seamGreedyEventuallyRight_iff_half_not_mem}{its transfer to non-membership};
\mword{Erdos249257/HalfCarryReachability.lean}{589}{finite_boolSupport_ne_half}{the finite-support exclusion};
\mword{Erdos249257/BooleanMobiusSkipRowCofinal.lean}{55}{exactLocalMersenneHalfRow_of_positiveHalfGreedySkip}{an exact finite sum from one greedy skip};
\mword{Erdos249257/BooleanMobiusSkippedCoreExactRow.lean}{228}{exists_upperHalfBooleanWord_of_skippedCore}{the upper-half Boolean fill};
\mword{Erdos249257/BooleanMobiusSkipRowCofinal.lean}{32}{greedyMersenneRemainderRat_half_pos}{the strict-positivity theorem};
\mword{Erdos249257/BooleanMobiusSkipRowCofinal.lean}{22}{CofinalPositiveHalfGreedySkips}{the cofinal-skip hypothesis};
\mword{Erdos249257/BooleanMobiusSkipRowCofinal.lean}{84}{cofinalExactLocalMersenneHalfRows_of_positiveHalfGreedySkips}{exact finite sums from arbitrarily late greedy skips};
\mword{Erdos249257/BooleanMobiusCofinalExactRows.lean}{71}{half_mem_mersenneAchievementSet_of_cofinalExactLocalRows}{compactness from exact sums at unbounded depths};
\mword{Erdos249257/BooleanMobiusSkipRowCofinal.lean}{97}{half_mem_mersenneAchievementSet_of_positiveHalfGreedySkips}{membership from infinitely many greedy skips};
\mword{Erdos249257/BooleanMobiusSkipRowCofinal.lean}{110}{cofinalPositiveHalfGreedySkips_iff_half_mem}{equivalence of infinitely many skips and half-membership};
\mword{Erdos249257/TerminalOnlyScaledVanishing.lean}{165}{half_mem_mersenneAchievementSet_of_terminalScaledVanishing}{achievement-set conclusion};
\mword{Erdos249257/TerminalOnlyScaledVanishing.lean}{221}{exists_infinite_support_half_of_terminalScaledVanishing}{infinite-support lift};
\lword{Erdos257/HalfCounterexampleFrontier.lean}{31}{exists_rational_half_counterexample_of_terminalScaledVanishing}{a conditional infinite support representing one half};
\lword{Erdos257/HalfCounterexampleFrontier.lean}{39}{not_universal_of_terminalScaledVanishing}{the conditional counterexample to universal irrationality}.

\paragraph{Finite greedy inequalities.}
\mword{Erdos249257/HalfGreedyTwoThirdsBand.lean}{88}{postTakeUnsafeAt_iff_band}{general band localization};
\mword{Erdos249257/HalfGreedyTwoThirdsBand.lean}{127}{singletonUnsafe_iff_twoThirdsBand}{two-thirds band};
\mword{Erdos249257/HalfGreedyTwoThirdsBand.lean}{231}{seven_le_of_intBand_odd}{odd numerator bound};
\mword{Erdos249257/HalfGreedyTwoThirdsBand.lean}{185}{not_twoThirdsBand_of_int}{integral safety}.

\paragraph{Fixed-core supports and dilation.}
\mword{Erdos249257/SupportSunflowerDichotomy.lean}{531}{sunflowerForcedCarrySupply_of_orthogonalPetalBouquet}{carry certificates for the stated fixed-core support};
\mword{Erdos249257/SupportSunflowerDichotomy.lean}{540}{irrational_erdosSupportSeries_of_orthogonalPetalBouquet}{irrationality for the stated fixed-core support};
\mword{Erdos249257/SupportSunflowerDichotomy.lean}{406}{SunflowerForcedSlotTailSelection}{the uniform tail-selection hypothesis};
\mword{Erdos249257/CompositeDilationDefect.lean}{30}{supportCoeff_mul_eq_add_defect}{the divisor-count identity under dilation};
\mword{Erdos249257/CompositeDilationDefect.lean}{103}{compositeDilationDefect_eq_zero_of_prime_support}{vanishing of the dilation correction for prime support};
\mword{Erdos249257/CompositeDilationDefect.lean}{119}{supportCoeff_mul_prime_support}{prime specialization};
\mword{Erdos249257/CompositeDilationDefect.lean}{133}{mem_compositeDilationDefect_bouquet_imp_exceptional_or_petal_dvd}{classification of the extra divisors};
\mword{Erdos249257/CompositeDilationDefect.lean}{151}{compositeDilationDefect_le_exceptional_add_petalCoeff}{the upper bound for the dilation correction};
\mword{Erdos249257/CompositeDilationDefect.lean}{218}{compositeDilationDefect_two_six_fixture}{the example with exponents two and six}.

\paragraph{The shared-prefix families.}
\mword{Erdos249257/SuffixCylinderTerminalOnlyBridge.lean}{264}{cofinalTerminalOnlyStrip_of_cofinalCylinderStages}{terminal bounds from the shared-prefix families};
\mword{Erdos249257/SuffixCylinderTerminalOnlyBridge.lean}{277}{exists_infinite_support_half_of_cofinalCylinderStages}{infinite half-support};
\mword{Erdos249257/SuffixCylinderTerminalOnlyBridge.lean}{287}{exists_infinite_positive_support_half_of_cofinalCylinderStages}{positive-support lift}.

\paragraph{The M\"obius example and two-prime digit systems.}
\mword{Erdos249257/MobiusSignSupportNoGo.lean}{111}{tsum_negativeMobius_eq_half_add_positiveMobiusTail}{identity};
\mword{Erdos249257/MobiusSignSupportNoGo.lean}{150}{half_add_one_div_sixty_three_le_tsum_negativeMobius}{bound};
\lword{Erdos257/HalfCounterexampleFrontier.lean}{61}{finiteErdosSum_ne_one_div_twenty_one}{exclusion of finite support for $1/21$};
\mword{Erdos249257/Primitive23Multiplicity.lean}{52}{exists_primitive23_solution_of_eleven_le}{existence for the stated two-prime digit system from rank eleven};
\mword{Erdos249257/Primitive23Multiplicity.lean}{26}{no_primitive23_solution_ten}{nonexistence for that system at rank ten};
\mword{Erdos249257/Primitive23Multiplicity.lean}{38}{exists_two_primitive23_solutions_eleven}{two solutions at rank eleven};
\mword{Erdos249257/Primitive23Multiplicity.lean}{86}{exists_two_primitive23_solutions_mul_ten}{two solutions at the stated multiples of ten}.

\paragraph{Scaled greedy remainders.}
\mword{Erdos249257/GreedyTrapDynamics.lean}{262}{mersenneAchievementSet_eq_scaledGreedyTrap}{membership and a bounded scaled greedy remainder};
\mword{Erdos249257/GreedyTrapDynamics.lean}{189}{scaledGreedyRemainder_tendsto_atTop_of_not_mem}{divergence of the scaled remainder outside the achievement set};
\mword{Erdos249257/GreedyTrapDynamics.lean}{225}{mem_mersenneAchievementSet_iff_scaledRemainder_cofinallyBounded}{membership and bounded scaled returns at unbounded ranks};
\mword{Erdos249257/GreedyTrapDynamics.lean}{152}{rat_mem_mersenneAchievementSet_iff_scaledLowerBranchCofinally}{the equivalent small-remainder test for rational targets};
\mword{Erdos249257/GreedyTrapDynamics.lean}{283}{one_div_twentyOne_mem_iff_scaledLowerBranchCofinally}{the small-remainder test for $1/21$};
\mword{Erdos249257/BooleanMobiusCarry.lean}{2117}{twentyOneGreedyDefect_add_mod_div_eq_scaled_remainder_add_tail}{the relation between the integer and scaled real remainders}.

\paragraph{Finite quotient tests for $1/21$.}
\mword{Erdos249257/TwentyOneQuotientGreedy.lean}{32}{twentyOneEvenQuotientGreedySupport}{support};
\mword{Erdos249257/TwentyOneQuotientGreedy.lean}{39}{twentyOneEvenQuotientGreedyRemainder}{remainder};
\mword{Erdos249257/TwentyOneQuotientGreedy.lean}{3507}{one_div_twenty_one_mem_iff_not_fatalAlignedBranch}{the equivalence with failure of $\mathcal F_{21}$};
\mword{Erdos249257/TwentyOneQuotientGreedy.lean}{5554}{twentyOneCofinalEvenQuotientGreedyDecay_of_closedRows}{vanishing scaled error from unbounded closed rows};
\mword{Erdos249257/TwentyOneQuotientGreedy.lean}{5658}{twentyOneFatalAlignedBranch_eventually_affine_supercapacity}{the eventual affine recurrence on $\mathcal F_{21}$};
\mword{Erdos249257/TwentyOneQuotientGreedy.lean}{231}{twentyOneClosedRow_forces_quotientGreedy}{uniqueness of a finite representation with the stated remainder bound};
\mword{Erdos249257/TwentyOneQuotientGreedy.lean}{5179}{twentyOneAlignedSaturatedCrossing_forces_canonical_ancestor_hole}{the missing-prefix consequence of an aligned crossing};
\mword{Erdos249257/TwentyOneQuotientGreedy.lean}{5223}{twentyOneAlignedSaturatedCrossing_forces_scaled_greedy_skip}{the real greedy skip forced by that crossing};
\mword{Erdos249257/GreedyTrapDynamics.lean}{275}{one_div_twentyOne_mem_iff_scaledRemainder_cofinallyBounded}{the bounded-return test for $1/21$};
\mword{Erdos249257/BooleanMobiusCarry.lean}{2843}{one_div_twenty_one_mem_mersenneAchievementSet_of_defect_le_one_cofinally}{a sufficient integer-remainder bound for $1/21$};
\mword{Erdos249257/BooleanMobiusCarry.lean}{1892}{twentyOneGreedyDefect_add_six}{the six-step integer-remainder recurrence};
\mword{Erdos249257/BooleanMobiusCarry.lean}{1927}{twentyOneGreedyDefect_add_six_le_div_six_iff_repairLoad}{the equivalent divisor-count inequality};
\mword{Erdos249257/TwentyOneQuotientGreedy.lean}{3735}{lastTwentyOneSkip_erdosBorwein_fatalInterval}{an exact one-sided approximation};
\mword{Erdos249257/TwentyOneQuotientGreedy.lean}{3569}{lastTwentyOneSkip_oddDoublingOrder_eq_lcm}{order identity};
\mword{Erdos249257/TwentyOneQuotientGreedy.lean}{3583}{lastTwentyOneSkip_le_oddDoublingOrder}{lower bound}.

\section*{Acknowledgements}
I thank Wouter van Doorn for advice on mathematical exposition, in
particular on explaining restrictive hypotheses, avoiding unnecessary
notation, and writing for a first-time reader.  His comments concerned a
note on Problem~243; this acknowledgement does not imply that he reviewed
the mathematics of the present paper.

\begin{thebibliography}{99}
\bibitem{duverneytachiya} D.~Duverney and Y.~Tachiya,
  \href{https://danielduverney.fr/documents/theorie-des-nombres/DuverneyTachiya190522.pdf}
  {\emph{Refinement of the Chowla--Erd\H{o}s method and linear independence
  of certain Lambert series}},
  Forum Math.\ 31 (2019), no.~6, 1557--1566,
  \href{https://doi.org/10.1515/forum-2018-0299}{DOI}.  Page numbers refer to
  the linked author preprint.
\bibitem{kanekosuzukitachiya} H.~Kaneko, Y.~Suzuki, and Y.~Tachiya,
  \href{https://arxiv.org/abs/2601.20743v1}{\emph{Refinements of Erd\H{o}s's
  irrationality criterion for certain sparse infinite series}},
  arXiv:2601.20743v1 (2026).
\bibitem{taoteravainen2025} T.~Tao and J.~Ter\"av\"ainen,
  \href{https://arxiv.org/abs/2512.01739v2}{\emph{Quantitative correlations and
  some problems on prime factors of consecutive integers}}, arXiv:2512.01739v2
  (submitted December 2025, revised April 2026).
\bibitem{kovactao} V.~Kova\v{c} and T.~Tao, \emph{On several irrationality
  problems for Ahmes series}, Acta Math.\ Hungar.\ 175 (2025), no.~2, 572--608,
  \href{https://doi.org/10.1007/s10474-025-01528-0}{DOI}.  Page numbers refer
  to arXiv:2406.17593v4.
\bibitem{erdos1948} P.~Erd\H{o}s, \emph{On arithmetical properties of Lambert
  series}, J.~Indian Math.\ Soc.\ 12 (1948), 63--66.
\bibitem{formalconjectures257} The Formal Conjectures Authors,
  \href{https://github.com/google-deepmind/formal-conjectures/blob/f776d2f2039351b00737ffcafb9d7d7666e1d9af/FormalConjectures/ErdosProblems/257.lean}
  {\emph{FormalConjectures.ErdosProblems.\texttt{257}}},
  Lean source at commit \texttt{f776d2f}, 2025, accessed 13 September 2026.

\bibitem{erdos1968} P.~Erd\H{o}s,
  \href{https://users.renyi.hu/~p_erdos/1969-09.pdf}{\emph{On the
  irrationality of certain series}}, Math.\ Student 36 (1968), 222--226
  (issued 1969); \href{https://www.renyi.hu/~p_erdos/1969-09.pdf}{five-page scan}.
\bibitem{erdos1975} P.~Erd\H{o}s, \emph{Some problems and results on the
  irrationality of the sum of infinite series}, J.~Math.\ Sci.\ 10 (1975),
  1--7.
\bibitem{bkkkz2026} K.~Barreto, J.~Kang, S.-H.~Kim, V.~Kova\v{c}, and
  S.~Zhang, \href{https://arxiv.org/abs/2601.21442v3}{\emph{Irrationality of
  rapidly converging series: a problem of Erd\H{o}s and Graham}},
  arXiv:2601.21442v3 (2026), to appear in Bull.\ London Math.\ Soc.
\bibitem{nitecki2013} Z.~Nitecki,
  \href{https://arxiv.org/abs/1106.3779v2}{\emph{Subsum sets: intervals, Cantor
  sets, and Cantorvals}}, arXiv:1106.3779v2 (2013).
\bibitem{lucatachiya2017} F.~Luca and Y.~Tachiya,
  \href{https://www.kurims.kyoto-u.ac.jp/~kyodo/kokyuroku/contents/pdf/2014-14.pdf}
  {\emph{Linear independence results for the values of divisor functions
  series}}, RIMS K\^oky\^uroku No.~2014 (2017), 138--150.
  Theorem~A restates their periodic-sequence theorem
  (\href{https://doi.org/10.1142/S1793042113501121}{DOI}).
\bibitem{lucatachiya2014periodic}
F.~Luca and Y.~Tachiya,
\emph{Irrationality of Lambert series associated with a periodic sequence},
International Journal of Number Theory \textbf{10} (2014), no.~3, 623--636.
\href{https://doi.org/10.1142/S1793042113501121}{doi:10.1142/S1793042113501121}.

\bibitem{lucatachiya2014independence}
F.~Luca and Y.~Tachiya,
\emph{Linear independence of certain Lambert series},
Proceedings of the American Mathematical Society \textbf{142} (2014), no.~10, 3411--3419.
\href{https://doi.org/10.1090/S0002-9939-2014-12102-2}{doi:10.1090/S0002-9939-2014-12102-2}.

\bibitem{vanassche2001}
W.~Van Assche,
\emph{Little $q$-Legendre polynomials and irrationality of certain Lambert series},
The Ramanujan Journal \textbf{5} (2001), 295--310.
\href{https://doi.org/10.1023/A:1012930828917}{doi:10.1023/A:1012930828917};
\href{https://arxiv.org/abs/math/0101187v1}{arXiv:math/0101187v1}.

\bibitem{hornich1941}
H.~Hornich,
\emph{\"Uber beliebige Teilsummen absolut konvergenter Reihen},
Monatshefte f\"ur Mathematik und Physik \textbf{49} (1941), 316--320.
\href{https://doi.org/10.1007/BF01707309}{doi:10.1007/BF01707309}.

\bibitem{endpoint2026} \emph{The logarithmic endpoint fails under arithmetic
  sampling}, AI-assisted ordinary proof note, 17 September 2026.
  Theorem~1, Corollary~3 and Proposition~4.  Unpublished working note;
  independent review and fresh Lean verification are outstanding.
\bibitem{vandoornkovac} W.~van Doorn and V.~Kova\v{c},
  \emph{Lacunary sequences whose reciprocal sums represent all rational
  numbers in an interval}, Acta Arith.\ 223 (2026), 275--295.
  \href{https://doi.org/10.4064/aa251001-13-1}{DOI}.
  Page references use \url{https://arxiv.org/abs/2509.24971v3}.

\end{thebibliography}

\companionpapercontext

\end{document}
